\section{A scale-dependent cutoff and entry}\label{sec:entry}

We construct a small local interaction for the critical Bessel heights
at a starting side that may increase with the size of the square.
Away from the magnetic source the construction gives the small
activities required by the height map.  At the source it gives one
labelled activity, whose norm may grow as a fixed power of the
observation-cell side.  The cutoff that makes these bounds possible
has a negligible cost in the finite physical partition functions.

\subsection{The averaged interaction and its subset targets}

From now on \(p=p_{b_c}\), \(\alpha=\sqrt a\), and
\(\omega=2\pi/\alpha\).  Let \(G\) be a finite retained graph whose
primary vertices are partitioned into whole square cells of side \(s\);
the estimates below concern the specified tile and box topologies.
Let \(B_s\) take the average in every retained cell.  Introduce
independent standard Gaussian noises \(e\) and observations
\[
                         y=B_sh+e.
\]
For prescribed \(y\), let \(Z_G^\tau(y)\) include the penalty
\(\exp(-\|B_sh-y\|^2/2)\), with every retained observation row.
The corresponding unpinned real Gaussian energy is
\begin{equation}\label{eq:observation-energy}
 E_G^\tau(y)=\min_u
 \left\{a^{-1}\norm{\nabla u+\tau}_G^2+\norm{B_su-y}^2\right\}.
\end{equation}
On a cut graph the minimum varies every free component constant.
At fixed \(y\), the observation penalty confines those constants in
the height sum as well; all component translation copies are summed.

We use the reference representation on the connected full graphs:
a compatible torus in the centered sector, or an aligned free square
in its centered or center-unit sector.  On such a graph choose \(v\)
so that \(\eta=\nabla v+\tau\) is co-closed, including at a free
boundary.  Orthogonality to gradients gives
\begin{equation}\label{eq:sector-completion}
 E_G^\tau(y)=a^{-1}\norm{\eta}_G^2+E_G^0(y-B_sv).
\end{equation}
Let \(N_{\rm c}\) be the number of observation cells and
\(\mathcal Y_G=\R^{N_{\rm c}}/(2\pi\Z\mathbf1)\).  Integrals on
this quotient use ambient Lebesgue measure on a fundamental domain.
Write \(Z_G^\tau\) without an observation argument for the bare
partition sum, modulo one common height translation.  Summing all
height translation copies in \(Z_G^\tau(y)\) and integrating their
Gaussian penalties gives exactly
\[
       \int_{\mathcal Y_G}Z_G^\tau(y)\,\dd y
               =(2\pi)^{N_{\rm c}/2}Z_G^\tau .
\]

Let \(\phi_0\) have covariance \(A^+\) on the mean-free modes,
let \(\theta\) be independent and uniform on \([0,\omega)\), and
let \(e\) be independent standard observation noise.  Put
\[
 Y_0=B_s\alpha(\phi_0+\theta\mathbf1)+e,\qquad
 Y_\tau=B_sv+Y_0 .
\]
Their quotient densities are
\[
 \frac{\dd\mu_G^0}{\dd y}
       =\kappa_Ge^{-E_G^0(y)/2},\qquad
 \frac{\dd\mu_G^\tau}{\dd y}
       =\kappa_G
          e^{-\{E_G^\tau(y)-a^{-1}\|\eta\|_G^2\}/2}.
\]
The same \(\kappa_G\) includes the common-mean Jacobian and cancels.
Consequently the exact reference identity is
\begin{equation}\label{eq:noisy-reference}
 \frac{Z_G^\tau}{Z_G^0}
 =e^{-\|\eta\|_G^2/(2a)}
 \frac{\E\!\left[e^{E_G^\tau(Y_\tau)/2}Z_G^\tau(Y_\tau)\right]}
      {\E\!\left[e^{E_G^0(Y_0)/2}Z_G^0(Y_0)\right]} .
\end{equation}
The same calculation applies when either observation integrand is
multiplied by a function on \(\mathcal Y_G\) for which the integrals
are finite, including the cutoff defined below.  This reference identity uses the connected full
graph and its co-closed \(v\).  The cut-graph sums used later instead
factor through their confined component observations.

Initialization will be used on compatible tori, in local prescriptions
on the plane, and on aligned free height squares.  The free height
grid for \(M_n\) has \(2n\) primary dual vertices on a side and one
circulation \(2\pi\) at its center.  Choose a straight primal current
path from the center to a side along observation-cell divisions, and
support \(\tau\) on the dual bonds crossing that path.  Use the same
cut direction and grid origins for every
square.  In particular the center is the junction of four observation
cells.  Let \(v_n\) be real vertex values such that
\[
       \eta_n=\nabla v_n+\tau,\qquad
       \ip{\eta_n}{\nabla f}=0\quad\hbox{for every real }f.
\]
The latter identity includes the free boundary and characterizes the
minimum-energy connection.  Write \(\eta_\infty\) and \(v_\infty\) for
the corresponding centered plane connection and branch.  Plane
expressions in this section are local prescriptions; there is no
infinite-plane sector partition integral.

Here is the exact Gaussian split used in the construction.  The
dyadic ratio \(L\) will be fixed in the parameter choice below.  With the
four nearest steps, the notation of \cite[Section~8]{B} has
\(Q=A/4\).  For an integral \(u\ge1\), put
\[
 P_u(\cos t)=
   \left(\frac{\sin(ut/2)}{u\sin(t/2)}\right)^{64},
 \qquad \mathsf P_u=P_u(1-A/64).
\]
The removable values are understood.  The split of
\cite[Equation~(9.11)]{B}, with no independent first piece, is
\begin{equation}\label{ent:reference-split}
\begin{split}
 C_{<u}&=A^{-1}(1-\mathsf P_u),\\
 \Gamma_u&=A^{-1}(\mathsf P_u-\mathsf P_{Lu}),\\
 C_{\ge u}&=A^+\mathsf P_u\quad\hbox{on nonconstant modes}.
\end{split}
\end{equation}
The first two quotients use their polynomial continuations at zero.
They are positive, have ranges \(O(u)\) and \(O(Lu)\), respectively,
and telescope on nonconstant modes.  Their continued constant modes
may be integrated along with the other modes: the final common mean is
uniform modulo \(\omega\), so adding an independent Gaussian constant
does not change a periodic integral.  On a free square these operators
are the even restrictions of the operators on the doubled torus.

Orient the retained faces between observation cells.  If a face
\(f=(i,j)\) is retained, the chosen connection is constant on its
microscopic crossing edges; call that value \(\tau_f\).  For the
observations of \eqref{eq:observation-energy}, set
\begin{equation}\label{ent:cutoff-definition}
\begin{split}
 w_f&=y_j-y_i+\tau_f,\\
 \chi_s(z)&=
 \frac{\displaystyle\int_{-s}^{s}\exp\{-(z-t)^2/2\}\dd t}
      {\displaystyle\int_{-s}^{s}\exp\{-t^2/2\}\dd t},\\
 \widehat Z_G^\tau(y)&=
          Z_G^\tau(y)\prod_{f\in\mathcal F(G)}\chi_s(w_f).
\end{split}
\end{equation}
Here \(\mathcal F(G)\) is the set of retained cell faces.  A cut
removes its cutoff factor along with its microscopic bonds.
The function \(\chi_s\) is entire and even; on the real line it lies
in \((0,1]\) and equals one at zero.  For \(s\ge1\), Gaussian tails
give
\begin{equation}\label{ent:cutoff-tails}
\begin{split}
 1-\chi_s(x)&\le Ce^{-cs^2}\quad (|x|\le s/2),\\
 \chi_s(x)&\le C\exp\{-\tfrac12(|x|-s)^2\}
                                  \quad (|x|\ge s).
\end{split}
\end{equation}

Fix the connection and collar radii in the block calculus of
\cite[Section~4]{B}.  At side \(u\), a polymer \(X\) is a finite
collection of blocks connected within that fixed block distance;
\(|X|\) denotes its block cardinality.  Let \(D_u(X)\) consist of
the sites in the blocks of \(X\) and in all grid blocks at index
distance at most two from \(X\), as in \cite[Section~4]{B}.
Difference stencils based there remain in a fixed microscopic
stencil neighborhood, denoted \(D_u(X)^+\) when needed.  At the
aligned scales used here, block sides are multiples of \(s\), and
the finite grids and clipped end blocks are whole \(s\)-cell arrays.
An open energy \(E_{D_u(X)}\) uses the corresponding union of whole
observation cells, all primary bonds with both endpoints in that union
(including intercell bonds), and every one of their observation rows; each resulting
component constant is free to minimize.  Compatible polymers have
separated collars.  For a
finite grid \(\mathcal B_u\), the subset functional is
\begin{equation}\label{eq:subset-functional}
 \cZ_u(K;\phi)=
 \sum_{V\subset\mathcal B_u}
       \prod_{X\in\Comp(V)}K(X,\phi),
\end{equation}
where components use the fixed connection radius and the empty
product is one.

At the compatible entry side \(m\) selected in
\eqref{ent:entry-side}, define for a graph \(\mathcal G\) in one of
the connected full classes just listed the interaction after entry averaging by
\begin{equation}\label{ent:averaged-interaction}
\begin{split}
 \mathcal H_{\mathcal G,m}^\tau(\psi)
 &=\E_{\zeta_{<m},e}
 \left[
   e^{E_{\mathcal G}^\tau(y)/2}
        \frac{\widehat Z_{\mathcal G}^\tau(y)}
             {Z_{\mathcal G}^0(0)}
 \right]_{\,y=B_s(v+\alpha(\psi+\zeta_{<m}))+e},\\
 &\hspace{18mm}\Cov(\zeta_{<m})=C_{<m}.
\end{split}
\end{equation}
For the centered sector \(v=0\); for a square unit sector \(v=v_n\).
The noises and the Gaussian field are independent.  Subsequent
integration of \(\mathcal H_{\mathcal G,m}^\tau\) and
\(\mathcal H_{\mathcal G,m}^0\) against the remaining covariance
\(C_{\ge m}\) and the uniform common mean gives the two expectations
in the cutoff version of \eqref{eq:noisy-reference}, divided by the
same \(Z_{\mathcal G}^0(0)\).  Their ratio therefore represents the
cutoff physical sector ratio after multiplication by
\(\exp\{-\|\eta\|^2/(2a)\}\), with \(\eta=\eta_n\) in the square
unit sector.  The remaining Gaussian covariance
is the same in both sectors.

The construction has an exact algebraic target before any norm is
chosen.  In the centered calculation it will produce a positive
field-independent scalar \(c_{\mathcal G,m}\) and ordinary component
activities \(K_{\mathcal G,m}^0\) satisfying
\begin{equation}\label{ent:ordinary-subset}
 \mathcal H_{\mathcal G,m}^0(\psi)
             =c_{\mathcal G,m}\cZ_m(K_{\mathcal G,m}^0;\psi).
\end{equation}
In the center-unit calculation it will use the same scalar and produce
unlabelled component activities \(K_{\mathcal G,m}^\tau\) together
with one labelled component \(J_{\mathcal G,m}\), so that
\begin{equation}\label{ent:labelled-subset}
 \mathcal H_{\mathcal G,m}^\tau(\psi)
   =c_{\mathcal G,m}[x]\,
       \cZ_m(K_{\mathcal G,m}^\tau+xJ_{\mathcal G,m};\psi).
\end{equation}
Here \([x]\) selects the coefficient of \(x\).  The component counted
by \(J\) contains the compulsory source marker and every dependency
attached to it; \(K^\tau\) counts the other components.  The common
scalar cancels between the two sectors of this square.  Theorem
\ref{ent:initialization} gives the parameter range, makes the ordinary
activities small in the norm defined next, and permits polynomial
growth of the labelled norm.  It also identifies the activity
prescriptions that copy exactly between aligned geometries.

To reach these identities with those bounds, we first obtain a
large-field reserve from the cutoff and isolate the single
circulation cost in the Gaussian energy.  The finite expansion then
records each component's complete occupied cover.  Its real estimate
passes from signed differences to positive halo topologies, cuts
those topologies into bounded pieces, and pays the resulting local
Gaussian costs.  An analytic strip estimate upgrades that real bound
to the derivative norm before the connected components are summed.

\subsection{The activity norm and the local map}

The local calculus below uses the exact partition map of
\cite[Proposition~4.8]{B}.  In a related analytic construction,
Bauerschmidt, Park, and Rodriguez begin their high-temperature discrete
Gaussian analysis with smoothing, then use the successive finite-range
Gaussian integrations described by Brydges \cite{Brydges2009} and a
localization of the noncontracting terms inspired by Brydges and Slade
\cite{BrydgesSlade2015}; see \cite[Section~1.5]{BPR2024I}.

The analytic point norm uses all real directional derivatives at a
fixed radius \(h_0\):
\begin{equation}\label{eq:point-norm}
 \begin{split}
 \norm{f}_{u,D}
 &=\max_{x\in D}\bigl\{|f(x)|,u|\nabla f(x)|,
                                      u^2|\nabla^2f(x)|\bigr\},\\
 |F|_{u,X,\phi}
 &=\sum_{k\ge0}\frac{h_0^k}{k!}
       \sup_{\norm{f_i}_{u,D_u(X)}\le1}
       |D^kF(\phi)[f_1,\ldots,f_k]|.
 \end{split}
\end{equation}
All fixed difference stencils are included in the padding.
The ordinary regulator \(W_u^\kappa(X,\phi)\) is the exponential
regulator in \cite[Equation~(4.10)]{B}: it controls the gradient
energy on \(D_u(X)\), its boundary trace, and scaled suprema of
second differences.  At a free side the field is extended evenly.
We use exactly that regulator, rather than a bound restricted to
bounded fields.

Use the accumulated covariance \(C_{<u}\) of
\eqref{ent:reference-split}.  Above a common absolute side, this
polynomial split is independent of the chosen entry side.
Finite-range decompositions by functional calculus have a general
antecedent in \cite{Bauerschmidt2013}; the exact polynomials and
boundary conventions here are those of \cite{B}.
For now regard \(p_1,p_2,h_*\) as parameters satisfying
\(1<p_1<p_2\) and \(0<h_*<1/p_2\).  They will be fixed with the
strict reserve supplied by \eqref{ent:large-field-reserve}.
With \(D=D_u(X)\), put
\begin{equation}\label{eq:observation-regulator}
 V_u(X,\phi)=\frac12\sup_{\zeta,e}
 \left\{h_* E_D^0\bigl(B_s\alpha(\phi+\zeta)+e\bigr)
 -p_1^{-1}\bigl(\|\zeta\|_{C_{<u}}^2+\|e\|^2\bigr)\right\},
\end{equation}
where \(\|\zeta\|_{C_{<u}}\) is the Cameron norm on the Gaussian
Hilbert space, including when the covariance is singular.  The full norm is
\begin{equation}\label{eq:activity-norm}
 \norm{F}_{u,A_w,V}
 =\sup_{X,\phi}
  A_w^{|X|}\frac{|F|_{u,X,\phi}}
                 {W_u^\kappa(X,\phi)e^{V_u(X,\phi)}}.
\end{equation}
When \(u=L^j\), we also write \(V_j=V_{L^j}\), with \(j\) the
absolute block index.
The block weight \(A_w\) is distinct from the Laplacian \(A\).

The following precise features of the cited calculus will be used.
They are analytic statements about activity functions and do not
assume that the functions came from a discrete Gaussian height model.

\begin{proposition}[Local map input]\label{prop:local-map}
Assume every activity has period \(\omega\) under common field translations,
and its field dependence lies in the sites of its occupied support
\(X\) together with a fixed microscopic collar of width at most a
fixed small fraction of \(u\).  The norm padding
\(D_u(X)\) is generally larger.  Call an input small when its connected
support has \(|X|\le n_{\mathrm{sm}}\), for the fixed block threshold
of the calculus.  For a linear direction assume the point bound
\[
 |F|_{u,X,\phi}\le
 \nu_X W_u^\kappa(X,\phi)e^{V_u(X,\phi)}
 \qquad\text{for every real }\phi,
\]
with \(\nu_X<\infty\); the resulting estimates are homogeneous in
\(\nu_X\).

For a sufficiently large fixed dyadic \(L\), sufficiently large
weights \(A_w\), and a sufficiently small activity ball, shell
integration and designated localizations on small linear inputs give
the exact identity
\begin{equation}\label{eq:exact-local-map}
 \E_\zeta\cZ_u(K;\phi+\zeta)
 =\prod_{B\in\mathcal B_{Lu}}(1+a_B)\,
        \cZ_{Lu}(K';\phi).
\end{equation}
The map is analytic in the activities.  If the total input size,
including separately measured singleton coefficients, is \(\nu\),
its nonlinear remainder is bounded by \(C_{L,A_w}\nu^2\), with the
corresponding differentiated bounds.  A residual from a large input retains its whole coarse
closure; only a piece from a small linear input can be assigned to a
singleton.  The coefficient \(a_B\) is the complete sum of the
field-independent pieces designated at \(B\).  Any undesignated
constant value of a translated nonconstant piece remains in \(K'\).

Write \(\Gamma_u=C_{<Lu}-C_{<u}\) for the shell from \(u\) to \(Lu\).
For its estimates on the aligned block and observation grids of a
finite square or torus of side \(M\), assume \(Lu\le M/D_0\), with
\(D_0\) a sufficiently large fixed range buffer.
The norms \eqref{eq:point-norm}--\eqref{eq:activity-norm} have the
shell composition, localization, charge, and large-support
contractions of \cite[Sections~4 and~8, Lemma~9.2]{B}.
In particular, the extra regulator is monotone under enlargement,
adds on compatible separated supports, and costs \(C^{|X|}\) under
shell averaging.  On the full bulk block grid, a small neutral input
invariant under field negation has gain \(CL^{-3}\) after constant and
quadratic localization.  Removing only the constant gives \(CL^{-2}\)
for a neutral input invariant under field negation and \(CL^{-1}\) for a general neutral
input, including at irregular end blocks.  For a charge-\(k\) component
of a small input, \(k\ne0\), the Gaussian charge gain is
\[
 \exp\left\{
 -\bigl(\beta_{\mathrm{sh}}|k|\alpha-\tfrac12\beta_{\mathrm{sh}}^2\bigr)
       \Gamma_u(v_0,v_0)+h_0|k|\alpha\right\}.
\]
Here a charge-\(k\) component obeys
\(F(\phi+c)=e^{ik\alpha c}F(\phi)\), \(v_0\) is a reference site
in a block of \(X\),
and \(\beta_{\mathrm{sh}}>0\) is the Cameron-shift amplitude.
The analytic radius must satisfy
\(h_0>C_{\mathrm{sh}}\beta_{\mathrm{sh}}\), where \(C_{\mathrm{sh}}\)
is the fixed geometric constant for that shift, and the parameters
are chosen so that
\(\beta_{\mathrm{sh}}\Gamma_u(v_0,v_0)>h_0\) by a uniform fixed
margin.  The displayed gains are then summable over \(k\ne0\).
This is \cite[Equation~(4.26)]{B}.  Large inputs can be made an
arbitrarily small linear contribution by increasing \(A_w\).
\end{proposition}

\begin{proof}
The exact identity and support assertions are
\cite[Proposition~4.8]{B}.  Its proof first groups compatible inputs
by coarse closure, then relocates chosen pieces from small linear
inputs, retaining their exclusion data, and finally multiplies by
\((1+a_B)^{-1}\) through constant decorations.  Thus it permits the
designated constants in the statement.  Its convergent multilinear
and scalar-inverse series give the analytic bounds.  The gains are
the localizations in \cite[Section~4]{B}.  The observation regulator
is \cite[Lemma~9.2]{B}; its fundamental estimate is contraction of
the open observation energy by the restricted Gaussian covariance.
The later free-wall, marked, and growing-terminal-grid uses require
the additional verifications given in
Sections~\ref{sec:entry} and~\ref{sec:observable}.
\end{proof}

On a regular block \(B\) of side \(u\), the two retained bulk
singletons are
\begin{equation}\label{eq:singletons}
 e_B^0(\phi)=\frac12\sum_{x\in B}\sum_{e=\pm e_1,\pm e_2}
                  |\nabla_e\phi(x)|^2,\qquad
 c_B^\alpha(\phi)=u^{-2}\sum_{x\in B}\cos(\alpha\phi_x).
\end{equation}
Summing \(e_B^0\) over the blocks gives the reference energy.
The weak map of \cite[Proposition~8.3]{B} retains
\(t e_B^0/2+z c_B^\alpha\); Section~\ref{sec:trajectory} states its
normal coordinates and the estimates used here.

\subsection{Reference and cutoff estimates}

The source background controls the deterministic cost of translating
the reference field.  The reflected covariance bounds and the
prescribed-observation cutoff bound then provide the reserves used at
entry.

\begin{lemma}[The centered unit background]\label{ent:center-background}
The branches can be chosen bounded, with a consistent additive
constant, and for \(0\le r\le4\) they satisfy, in lattice units,
\begin{equation}\label{ent:unit-decay}
\begin{split}
 |\nabla^r\eta_n(e)|
   &\le C_r(1+\dist(e,0))^{-1-r},\\
 |\nabla^r(\eta_n-\eta_\infty)(e)|
   &\le C_r n^{-1-r}
       \qquad\bigl(\dist(e,0)\le c_0n\bigr),
\end{split}
\end{equation}
where \(c_0>0\) is fixed and sufficiently small.  The first bound also
holds for \(\eta_\infty\).  Differences at a free side use the
reflected stencils of the local calculus.  On a source-free rectangle
with its required collar, there is a single-valued real lift
\(\lambda\) in reference-field units satisfying
\[
                         \nabla\lambda=\eta_n/\alpha .
\]
Changing the integral gauge changes this lift by an integer multiple
of the common period \(\omega=2\pi/\alpha\).
\end{lemma}

\begin{proof}
The connection is \(2\pi\) times the rotated gradient of the primal
Dirichlet Green function with pole at the center.  Rotation makes it
co-closed on the free dual square and gives the prescribed circulation.
The differentiated Green estimates of \cite[Lemma~4.2]{B} give the
first bound.  For the second, use the binary finite-range decomposition
in the proof of that lemma.  In the indicated central region all
pieces below a fixed multiple of \(n\) agree exactly with their plane
pieces.  A piece of side \(t\) contributes at most \(C_rt^{-1-r}\)
to the differentiated rotated gradient; the remaining dyadic tail
is at most \(C_rn^{-1-r}\).  Odd reflection of the primal potential
gives the stated free reflected convention for its rotated lift.

The circulation vanishes on a source-free rectangle, so path
integration there gives \(\lambda\); the integral gauge accounts for
its period ambiguity.  A bounded branch for one centered source can
be seen directly from the first bound.  By reflection, choose its
value constant on a reflection ray, or centered between the two rows
adjacent to that ray.  On a square ring of radius \(t\), the number
of increments is \(O(t)\), while their sizes are \(O((1+t)^{-1})\).
Integration from the ray around each ring is therefore bounded
uniformly in \(t\).  The finitely many increments at the source are
bounded as well.  Use the same ray and additive convention in the
square and in the plane.
\end{proof}

\begin{lemma}[Reflected and accumulated regulator bounds]\label{ent:entry-regulator}
At these compatible block and observation sides, with strict reserve in the parameters of
\eqref{eq:observation-regulator}, the shell composition,
localization, and charge estimates in Proposition~\ref{prop:local-map}
apply to \(\Gamma_u\) in \eqref{ent:reference-split} on an aligned
free square of side \(M\) when \(Lu\le M/D_0\), with \(D_0\) fixed
and sufficiently large, using only internal faces in every open
energy.  The accumulated covariance \(C_{<u}\) has the finite-range
and open-observation contraction bounds used in \(V_u\) and the
localized entry determinant.
\end{lemma}

\begin{proof}
The reference part of this assertion is the proof of
\cite[Lemma~9.2 and Lemma~4.1]{B}.  Restricting one trial field to all
components of an open union can only lower its minimized observation
energy; each component may minimize its constant separately.  Thus the
map from the Gaussian Hilbert coordinates of any accumulated pieces
and the noise into that energy has operator norm at most one.  This
is the covariance contraction used to optimize \(V_u\).  The same
trial restriction uses just internal faces at a free side.  Under
\(Lu\le M/D_0\), even reflection gives the positive-order bounds for
\(\Gamma_u\) there, so the trace,
H\"older, and small-support Taylor arguments in that proof apply with
the same strict parameter margins.
\end{proof}

\begin{lemma}[Physical and prescribed-observation bounds]\label{ent:cutoff-bounds}
There are \(c>0\) and \(s_1<\infty\), independent of the finite
cell topology, such that for every \(s\ge s_1\), every cut topology
used below, and its centered sector or the restriction of the aligned
unit sector,
\begin{equation}\label{ent:pointwise-cutoff}
 \frac{\widehat Z_G^\tau(y)}{Z_G^0(0)}
       \le \exp\left\{-c\sum_{f\in\mathcal F(G)}w_f^2\right\}
                    \qquad(y\hbox{ real}).
\end{equation}

For a connected full finite graph, let
\(\mathcal P_G^\tau(F)\) denote the integral of \(F(y)Z_G^\tau(y)\)
on the observation space modulo \(y\mapsto y+2\pi k\mathbf1\),
\(k\in\mathbb Z\), and let \(\widehat{\mathcal P}_G^\tau(F)\) use
\(\widehat Z_G^\tau\).  For the centered sector take either a uniformly
bounded function \(F\) on this quotient, or
\(F=\ell(y)^k e^{t\ell(y)}\) times a uniformly bounded function on
the quotient, where \(k\) is a fixed nonnegative integer,
\(t\) is in a fixed compact subset of \(\R\), and the mean-free linear observation
\(\ell\) has uniformly bounded primary reference cost and noise
variance.  For the square unit sector take \(F=1\).
Then
\begin{equation}\label{ent:cutoff-loss}
 \frac{\left|\mathcal P_G^\tau(F)
             -\widehat{\mathcal P}_G^\tau(F)\right|}
      {\mathcal P_G^0(1)}
       \le C_F|\mathcal F(G)|e^{-cs^2}.
\end{equation}
In particular, if the logarithm of the side is \(O(N)\) and
\(s\asymp N\), this bound is at most \(e^{-c'N^2}\) for some
\(c'>0\).  The same conclusion holds for the difference of normalized
centered averages of the stated tests.
\end{lemma}

\begin{proof}
For adjacent cells \(i,j\), transport each site in \(i\) straight
along its row to its translate in \(j\), and average these paths.
This writes the difference of the two cell averages, with its
connection value, as
\[
       (B_sh)_j-(B_sh)_i+\tau_f
              =\sum_e\Phi_f(e)(\nabla_eh+\tau_e).
\]
The flow \(\Phi_f\) is supported on the two cells, has coefficients
at most \(C/s\), and has squared coefficient sum at most \(C\).
Indeed each of the \(s^2\) paths has weight \(s^{-2}\) and length
\(s\), and any primary edge belongs to at most \(s\) of them.
There is bounded overlap among the flows for incident faces.
Consequently
\begin{equation}\label{ent:cell-flow}
 \max_e\left|\sum_f t_f\Phi_f(e)\right|
       \le Cs^{-1}\max_f|t_f|,
 \qquad
 \sum_e\left|\sum_f t_f\Phi_f(e)\right|^2
       \le C\sum_f t_f^2.
\end{equation}
The noise part \(\sum_ft_f(e_j-e_i)\) has variance at most
\(C\sum_ft_f^2\).

Fix any \(k_0<\infty\).  The tilted affine comparison of
Proposition~\ref{prop:comparisons}, followed by release to independent
primary increments, therefore gives
\begin{equation}\label{ent:face-tilt}
 e^{\sum_ft_fw_f}\frac{Z_G^\tau(y)}{Z_G^0(0)}
       \le \exp\{C_{k_0}\sum_ft_f^2\},
                 \qquad \max_f|t_f|\le k_0s .
\end{equation}
Here is how to apply the comparison at a prescribed continuous
observation.  On the chain subdivision in the proof of that
proposition use the height and noise coordinates \((h,e)\), with
the affine equality \(e=y-B_sh\).  The statistic before releasing
this equality is
\(\sum_f t_f[(B_sh)_j-(B_sh)_i+\tau_f+e_j-e_i]\); its value on the
equality is \(\sum_ft_fw_f\).  Release the equalities, then the
endpoint, curl, and period constraints, always with centered
division and this same affine statistic.  Constants may first be
given a vanishing auxiliary mass; the statistic annihilates them
and their Gaussian copy factors cancel.  The independent primary
transform is
\(\exp\{b_c(\cosh(2\pi x)-1)\}\).  By \eqref{ent:cell-flow} its
arguments are in a fixed bounded interval, where its logarithm is
at most \(C x^2\).  The independent noise gives the stated
quadratic cost.  Passing to the limit at the fixed primary graph
proves \eqref{ent:face-tilt}.  This argument does not use the bare
Gaussian variance of the subdivided chains.

Choose \(t_f=\delta\,\operatorname{sgn}(w_f)\min(|w_f|,s)\), with
\(\delta>0\) small compared with the constant in
\eqref{ent:face-tilt}.  Its exponent gives a negative multiple of
\(w_f^2\) for \(|w_f|\le2s\), and a negative multiple of
\(s|w_f|\) beyond that range.  On the latter faces use the second
bound in \eqref{ent:cutoff-tails}, for which
\((|w_f|-s)^2\ge w_f^2/4\).  The bounded prefactors are absorbed
by the \(s|w_f|\) gain once \(s\ge s_1\).  On all other faces
\(\chi_s\le1\).  Multiplication gives
\eqref{ent:pointwise-cutoff}.

The same release argument before prescribing observations gives
the integral version of \eqref{ent:face-tilt}, with
\(\mathcal P_G^0(1)\) as denominator.  Chernoff's inequality with
\(|t_f|\) a small multiple of \(s\) bounds the centered-normalized
integral of \(1_{\{|w_f|>s/2\}}\) by \(Ce^{-cs^2}\), in either
sector.  On its complement the first bound in
\eqref{ent:cutoff-tails} applies.  Since
\[
             1-\prod_f\chi_s(w_f)\le\sum_f(1-\chi_s(w_f)),
\]
this proves \eqref{ent:cutoff-loss} for bounded tests and for the
unit partition integral.  For a centered fixed Laplace tilt and its
fixed polynomial factors,
H\"older and the reference-cost exponential bound of
Proposition~\ref{prop:comparisons} give the same estimate with a
smaller \(c\); the noise variance is included in its cost.
The number of faces in a side-\(M\) graph is \(O((M/s)^2)\).
Thus \(\log M=O(N)\) and \(s\asymp N\) turn the right side into
\(e^{-c'N^2}\).  Finally the centered cutoff partition integral
is \(1+O(|\mathcal F(G)|e^{-cs^2})\) times its uncut value.
Dividing the two centered averages proves the last assertion.
\end{proof}

\subsection{Localizing the energy in face variables}

The Gaussian proxy in the isolation expansion must still be local
when its face values carry the unit circulation.  A comparison
matrix written only in the values \(y_i\) does not expose this point.
We therefore extend the observation minimum to independent real
face data.

The entry family in this subsection consists of the free retained
graphs produced by the finite tile isolation and box cuts below,
together with their current patch restrictions.  Every present
observation face is wholly retained or wholly cut, and every present
observation row is retained.  Natural clipped walls and lifted
periodic charts are allowed as in that construction.  A current
patch inherits all earlier cuts.  Temporary hard primary guards
belong only to the finite comparisons and are not members of this
free-energy family.

Regard the primary sites in each cell as a separate set, joining
them only across retained faces.  If \(w\in\R^{\mathcal F(G)}\),
write \(\nabla_G^w\xi\) for the ordinary difference inside a cell
and for \(\xi_j-\xi_i+w_f\) on a microscopic edge crossing the
oriented face \(f=(i,j)\).  Define
\[
 \mathcal Q_G(w)=\min_\xi
 \left\{a^{-1}\|\nabla_G^w\xi\|^2+\sum_i|(B_s\xi)_i|^2\right\}.
\]
For physical \(w_f=y_j-y_i+\tau_f\), the substitution
\(\xi|_i=u|_i-y_i\) proves
\(\mathcal Q_G(w)=E_G^\tau(y)\).  If all four faces incident to a
cell vertex \(z\) are retained, let
\(\sigma_z(w)\) be their oriented sum around \(z\).  There is no
such variable at a junction with a cut face.  Denote the set of
full junctions by \(\mathcal V_4(G)\).

\begin{lemma}[Face energy and changes at cuts]\label{ent:face-localization}
For the entry family just described, there is a number \(D_s\),
the same at every full junction,
and a symmetric matrix \(H_G\) such that
\begin{equation}\label{ent:face-decomposition}
 \mathcal Q_G(w)=D_s\sum_{z\in\mathcal V_4(G)}\sigma_z(w)^2
             +\sum_{f,f'\in\mathcal F(G)}H_G(f,f')w_fw_{f'},
 \qquad 0\le D_s\le C\log(2+s).
\end{equation}
The second matrix need not be positive.  In cell distances it has
the bounds
\begin{equation}\label{ent:face-matrix-bounds}
\begin{split}
 |H_G(f,f')|&\le Ce^{-c\,d(f,f')},\\
 |H_G(f,f')-H_{\widetilde G}(f,f')|
 &\le C\min\left\{
       e^{-c\,d(f,f')},
       e^{-c\,d(\{f,f'\},\mathcal D)}
                    \right\}.
\end{split}
\end{equation}
In the second line \(\mathcal D\) contains the cells and faces where
the two topologies or their local boundary rules differ, enlarged by
the fixed stencil radius of the local trial, forcing, and constant
quadratic coefficients used below.  Extend absent entries by zero
on the oriented face set of one common uncut inclusion \(U\), and
measure all distances there.  The two decays can be combined after
decreasing \(c\).

For the guarded comparison, let \(E\) be the same still-retained
local deletion in the current full graph and its current patch:
\[
 G_1=G_0\setminus E,\qquad
 P_0=G_0|_O,\qquad P_1=G_1|_O=P_0\setminus E.
\]
These four graphs belong to the stated tile and box family.  The
patch restrictions inherit every earlier cut and every retained
observation row and local rule in \(O\); their only new free wall is at the artificial boundary
of \(O\).  On the common face set define
\[
 \Delta=(H_{G_0}-H_{G_1})-(H_{P_0}-H_{P_1}).
\]
Let \(K_\Delta\) contain both before/after discrepancy sets and let
\(A_\Delta\) contain both full/patch discrepancy sets, including
all absent exterior cells and faces and the artificial-boundary
collar.  Enlarge both by the same fixed stencil radius.  If
\(d_U(K_\Delta,A_\Delta)\ge r\), then
\begin{equation}\label{ent:guarded-face-error}
 |\Delta(f,f')|
 \le C e^{-cr-c\,d_U(f,f')-c\,d_U(\{f,f'\},K_\Delta)}.
\end{equation}
Shared physical free walls and earlier cuts are common data and do
not belong to \(A_\Delta\).  Temporary hard primary guards are not
among these four free-energy graphs.
All constants are independent of \(s\) and of the containing
finite volume.  For the physical centered or unit data this also
gives
\begin{equation}\label{ent:face-energy-upper}
 E_G^\tau(y)\le C\sum_{f\in\mathcal F(G)}w_f^2
       +C\,1_{\{\text{unit junction is full in }G\}}\log(2+s).
\end{equation}
\end{lemma}

\begin{proof}
We first give a fixed local trial field linear in \(w\).  Use cell
units and one rooted coordinate and orientation convention, also in
lifted periodic charts.  Fix zero plateaus in the middle of the
cells, one local cell and quadrant triangulation, and nested junction
squares with half-side radii
\[
                  0<\rho_0<\rho_1<\rho_2<1/4.
\]
On an oriented retained face \(f:i\to j\), prescribe the two central
face traces \(+w_f/2\) in cell \(i\) and \(-w_f/2\) in cell \(j\).
Their compensated jump is zero.  A cut side has an independent
zero trace.

At a junction, the retained incidence mask is a subgraph of the
four-cycle.  If the junction is not full, this graph is a forest.
On each connected sector choose the unique mean-zero potential
\(c_z\) with \(dc_z=w\) on its retained oriented faces, and set
the inner cell values to \(-c_z\).  For an isolated cell set
\(c_z=0\).  These values have \(d\xi+w=0\) on the inner retained
edges.

At a full junction put \(\sigma=\sigma_z(w)\), orient its faces
cyclically, and write
\[
       w=w^{\rm grad}+w^{\rm circ},\qquad
       w_e^{\rm circ}=\sigma/4,\qquad dc_z=w^{\rm grad},
\]
with \(c_z\) again mean-zero.  Normalize the fixed plane angular
branch of Lemma~\ref{ent:center-background} to circulation one:
\[
       \eta^{\rm ang}=\nabla v^{\rm ang}+\tau^{\rm ang},
       \qquad \sum_{\rm cycle}\tau_e^{\rm ang}=1.
\]
Use the same ray and additive convention at every junction.
The four spoke jumps
\(\tau^{\rm ang}-(1/4,1/4,1/4,1/4)\) have zero cycle sum.
Let \(k_z\) be the mean-zero quadrant potential satisfying
\[
       dk_z=\tau^{\rm ang}-(1/4,1/4,1/4,1/4).
\]
The inner
field in quadrant \(i\) is
\[
 \xi^{\rm in}(x)=-c_z(i)
                +\sigma\bigl(v^{\rm ang}(x)+k_z(i)\bigr).
\]
It satisfies the exact compensated identity
\[
                 \nabla_G^w\xi^{\rm in}=\sigma\eta^{\rm ang}
\]
on the inner edges.  This identity fixes the sign of the chart
changes.  All choices are real and linear in the incident \(w\).

In each quadrant ring, use the fixed affine triangulation to extend
its two constant spoke traces, putting their mean at the inner and
outer diagonal nodes.  Call this bounded outer field \(\xi^{\rm out}\),
and retain \(\xi^{\rm in}\) for the inner field.  Blend them with one geometric cutoff \(\chi\)
on both sides of every retained spoke, equal to one through
\(\rho_1\) and zero at \(\rho_2\).  For an oriented microscopic
edge \(x\to y\), with \(w_e=0\) inside a cell and \(w_e=w_f\)
across a retained face, the exact product rule is
\[
\begin{split}
 d[\chi\xi^{\rm in}+(1-\chi)\xi^{\rm out}]_e+w_e
 ={}&\chi_x(d\xi^{\rm in}_e+w_e)\\
 &{}+(1-\chi_x)(d\xi^{\rm out}_e+w_e)\\
 &{}+(d\chi)_e(\xi^{\rm in}_y-\xi^{\rm out}_y).
\end{split}
\]
On this fixed ring both compensated differences on the right are
\(O(s^{-1}\sum|w_f|)\), \(|d\chi|\le C/s\), and
\(\xi^{\rm in}-\xi^{\rm out}\) is a
bounded linear combination of the incident face values.  Thus all
ring increments have the same \(O(s^{-1}\sum|w_f|)\) bound,
including the crossing edges.  The remaining cell triangulation
joins the outer traces and the central face traces to the zero
plateau.  Since \(\rho_2<1/4\), neighboring endpoint rings are
disjoint, and their intervening face segment uses the same prescribed
trace.  The cell averages are bounded linear combinations of the
local face values.  Sampling these prescriptions on the lattice
gives the trial field on every mask in the entry family, including
clipped masks.  A fixed dihedral average may enforce square symmetry
while retaining the common inner convention.

Use the complete symmetric edge set of the inner square of half-side
\(\rho_0\), with one fixed boundary-edge convention, and keep the
pure angular profile through \(\rho_1\).  Put
\[
                 D_s=a^{-1}\|\eta^{\rm ang}\|_{\rm inner}^2 .
\]
The plane bound in Lemma~\ref{ent:center-background} gives
\(0\le D_s\le C\log(2+s)\), and the chosen inner set makes \(D_s\)
identical at every full junction.  The gradient part has zero
compensated increment on that set.  Its complement in the pure
angular region is at distance at least \(cs\) from the junction.
The remaining trial increments
are \(O(s^{-1})\) times local face values on \(O(s^2)\) edges per
fixed stencil, and the cell averages are bounded.  Hence all their
quadratic coefficients are bounded and have finite cell range.
The trial already proves
\[
 \mathcal Q_G(w)\le
 D_s\sum_z\sigma_z(w)^2+C\sum_f w_f^2.
\]

We also need locality after minimization.  The quadratic form for
a correction to the trial is
\[
                 T_G=A_G/a+B_s^*B_s .
\]
It is coercive with constants independent of refinement.  In fact
cell Poincar\'e gives, for every complete cell \(Q\),
\begin{equation}\label{ent:cell-coercivity}
 s^{-2}\sum_{x\in Q}|\xi_x|^2
 \le C\left\{a^{-1}\sum_{e\subset Q}|\nabla_e\xi|^2
                         +|(B_s\xi)_Q|^2\right\}.
\end{equation}
Write \(u_G^{\rm tr}(w)\) for the trial just constructed and
\(C_G(w)\) for its energy.  Expanding that energy after adding a
correction \(\xi\) gives
\[
       \ip{\xi}{T_G\xi}+2F_G(w)(\xi)+C_G(w).
\]
Define its edge and observation coefficient arrays by
\[
 r_{G,f}=\partial_{w_f}\bigl(\nabla_G^w u_G^{\rm tr}(w)\bigr),
 \qquad
 b_{G,f}=B_s\partial_{w_f}u_G^{\rm tr}(w).
\]
Then \(F_G(w)=\sum_fw_fF_{G,f}\), with
\[
 F_{G,f}(\xi)=a^{-1}\sum_e r_{G,f}(e)\nabla_e\xi
                 +\sum_Q b_{G,f}(Q)(B_s\xi)_Q .
\]
A face affects only its two endpoint templates and the incident
cell interpolations.  Include this fixed decision neighborhood in
the forcing and discrepancy stencil.  Whenever the rooted retained
mask, alignment, and local boundary rule agree there, the arrays
\(r_{G,f},b_{G,f}\) agree exactly.  Each functional has that fixed
cell stencil and satisfies
\[
              |F_{G,f}(\xi)|\le C\|\xi\|_{T_G,\text{stencil}}.
\]
For the inner angular part this bound uses co-closure, and this
use is essential.  Choose a cutoff equal to one on its inner square
and supported inside the angular prescription.  The pairing of
the compensated angular gradient with the gradient of the cutoff
test vanishes by discrete co-closure.  The remaining product terms
lie on the ring, where the angular increments and cutoff
differences are \(O(s^{-1})\).  Cauchy--Schwarz and
\eqref{ent:cell-coercivity} bound them by the displayed local form
norm.  The observation residuals obey the same bound.  Thus the
logarithmic inner energy has left no logarithm in the dual norm of
the forcing.  Likewise
\(C_G(w)-D_s\sum_z\sigma_z(w)^2\) has bounded finite-range
coefficients.

Minimization now gives
\[
 \mathcal Q_G(w)
       =C_G(w)-\ip{F_G(w)}{T_G^{-1}F_G(w)}.
\]
The inverse in this formula has exponential energy locality for
these bounded local functionals.  To check this directly, let
\(T_Gu=F_{G,f}\) and test the equation away from its stencil by
\(\chi^2u\), where \(\chi\) changes over a fixed number of cells
and \(|\nabla\chi|\le C/s\).  Discrete product differences and
Cauchy--Schwarz absorb the outer gradient energy.  The transition
terms are bounded by \(s^{-2}\sum|u|^2\) on a fixed enlargement
of the layer, hence by its local energy by
\eqref{ent:cell-coercivity}.  Cells in which \(\chi\) is constant
give a nonnegative observation term; the observation cross terms
in transition cells obey the same local bound.  The outer energy
is consequently at most a fixed constant times the layer energy.
Increasing the layer by a fixed width reduces the remaining tail
by a fixed factor less than one.  Iteration proves exponential
decay.  This is the cell-cutoff proof underlying
\cite[Proposition~3.14]{B}, now applied to the local face
functionals just constructed.  It uses only retained bonds and
whole cells, so it is valid up to a free side.

For completeness, the same argument compares two topologies.
Suppose they agree out to cell distance \(D\) from a forcing
stencil.  Cut its solution to zero at distance of order \(D/2\)
inside the agreeing region.  Its change in the first form norm
is \(Ce^{-cD}\).  In the second topology its residual equation
has the same bound in dual norm: test after multiplication by
a second cutoff equal to one on the needed stencils and supported
in the agreeing region.  Multiplication by that cutoff is bounded
in local form norm by \eqref{ent:cell-coercivity}.  Coercivity in
the second topology bounds the difference from its solution by
\(Ce^{-cD}\).  Pair with another local forcing in the agreeing
region.  This proves the cut-distance part of
\eqref{ent:face-matrix-bounds}; direct solution-tail decay proves
its face-distance part.  Entries next to a changed face are
bounded directly.  Taking the geometric mean combines the two
bounds with smaller exponents.

For the four graphs in the guarded comparison, put
\[
 a=d_U(\{f,f'\},K_\Delta),\qquad
 b=d_U(f,f'),\qquad h=d_U(\{f,f'\},A_\Delta).
\]
Apply \eqref{ent:face-matrix-bounds} first to the two before/after
differences and then to the two full/patch differences.  Together
with the pair-distance bound, these give the three global estimates
\[
       |\Delta|\le Ce^{-ca},\qquad
       |\Delta|\le Ce^{-cb},\qquad
       |\Delta|\le Ce^{-ch}.
\]
These statements include absent entries.  If a face is absent from
a patch, it lies in \(A_\Delta\), so \(h=0\); the pair estimate
still controls a pair crossing the patch boundary.  Choose from
\(\{f,f'\}\) an endpoint nearest \(K_\Delta\) and one nearest
\(A_\Delta\).  The triangle inequality gives \(r\le a+b+h\), hence
\(2a+2b+h\ge r+a+b\).  The weighted geometric mean with weights
\(2/5,2/5,1/5\) now yields
\[
 |\Delta|\le C e^{-c(2a+2b+h)/5}
          \le C e^{-c'(r+a+b)},
\]
which is \eqref{ent:guarded-face-error} after decreasing \(c\).
Disconnected inclusion components are direct sums.  The retained
core-distance decay makes summation over deletion positions
harmless, since cell balls have polynomial volume.

For physical face data, the oriented sum of the \(y\)-differences
around a full junction is zero.  The connection has zero
circulation except for \(2\pi\) at the unit.  Inserting this in
the trial upper bound proves \eqref{ent:face-energy-upper}.
\end{proof}

The separated term in \eqref{ent:face-decomposition} is
deterministic in every Gaussian integration here.  Indeed, if
\(y=y^0+B_sv_n\), then
\begin{equation}\label{ent:fixed-circulation}
 \sigma_z\bigl(\nabla_{\rm c}y^0+\nabla_{\rm c}B_sv_n+\tau\bigr)
                 =\sigma_z(\tau),
\end{equation}
where \((\nabla_{\rm c}y)_f=y_j-y_i\).  It is zero or the fixed
number \(2\pi\).  Tile deletions preserve the source junction, so
this term cancels exactly in their before-minus-after proxy
differences.  An auxiliary seam used for a bound can remove the
junction, at a cost \(O(\log(2+s))\) once at the source.  In
\eqref{ent:fixed-circulation} neither the high Gaussian field nor
the observation noise is present on the right.  The term therefore
has zero Hessian in their coordinates and creates no Gaussian
determinant.  Imaginary field translations have gradient face data
as well, so the same observation applies to them.

\subsection{The initialization statement}

Combining \eqref{ent:pointwise-cutoff} with
\eqref{ent:face-energy-upper}, and decreasing a fixed \(c_*>0\) if
necessary, gives the large-field reserve
\begin{equation}\label{ent:large-field-reserve}
 e^{E_G^\tau(y)/2}\frac{\widehat Z_G^\tau(y)}{Z_G^0(0)}
 \le
 \exp\left\{\frac{1-c_*}{2}E_G^\tau(y)
       +C\,1_{\{\text{unit junction is full in }G\}}\log(2+s)\right\}.
\end{equation}
Indeed the upper bound for \(E_G^\tau\) bounds a fixed positive
multiple of that energy by the face-square sum plus the displayed
unit term.  We may and do take \(0<c_*<1\).  This is the reserve
that replaces the bare discrete-Gaussian large-field estimate in
the entry proof of \cite{B}.

We specify the parameter order.  First fix the interaction, connection,
and collar radii of the local calculus, and the ordinary regulator
parameters with the exponent reserves needed for H\"older.  Choose
\begin{equation}\label{ent:reserve-order}
       1-c_*<h_1<h_*<1/p_2,\qquad 1<p_1<p_2,
\end{equation}
leaving room to increase \(h_*\) and the ordinary regulator exponent
slightly.  These choices are possible because \(c_*>0\), and they
are made before \(R\) or \(s\).  Fix the analytic derivative radius
and imaginary strip widths needed for the marginal fundamental
frequency \(\alpha=\sqrt{8\pi}\), then a sufficiently large fixed
dyadic \(L\), then sufficiently large block weights \(A_w\), and
finally a sufficiently small activity ball, in the order of
\cite[Sections~4 and~8]{B}.  We may fix a finite number of such
weights and regulators at once.  In particular we keep an ordinary
unmarked weight stronger than the weight used after translating
activities by a source lift.

Next fix a dyadic coloring period \(q\) as in the finite isolation
expansion of \cite[Section~9]{B}.  Choose \(q\) larger than a fixed
multiple of the physical proximity, core, open-proxy, guard, testing,
and skip radii in units of \(R\) cells.  These dimensionless radii are
fixed before \(q\) and describe only the local isolation comparisons.
With \(q\) fixed, choose a large dyadic \(R\) and put \(P=R^5\).
The period \(q\) does not increase with \(R\) or \(s\).  The entry
side is
\begin{equation}\label{ent:entry-side}
                     m=qRs=L^{j_0}.
\end{equation}
The spatial quantities have the following distinct roles:
\begin{center}
\begin{tabular}{@{}c@{\quad}p{0.76\linewidth}@{}}
\(s\) & side of one observation cell, in primary lattice units;\\
\(R\) & fixed local isolation scale, in observation cells;\\
\(P=R^5\) & side of an auxiliary comparison box, in observation cells;\\
\(q\) & fixed coloring period, measured in units of \(R\) cells;\\
\(m\) & entry block side, in primary lattice units;\\
\(j_0\) & absolute block index of that entry side.
\end{tabular}
\end{center}
After \(R,P,q\) are fixed we choose a finite large-field threshold
\(U_{\rm big}\), and only then
take \(s\) sufficiently large among the dyadic values satisfying
\eqref{ent:entry-side}.  Such values are arbitrarily large.

The occupied support in the statement records all evaluated
coordinates and local rule data, together with the complete paths
that tag a link's winding or wall dependence.  We use its norm
padding \(D_m(X)\) as a sufficient region for the entry copy tests.
For a torus, say that \(D_m(X)\) \emph{unwraps} when its sites, bonds,
fixed stencils, and occupied path incidences admit one injective lift
to the aligned plane block grid.  A square copy test requires that
this region, including the fixed stencil neighborhood, miss the
wall.  For a square-to-plane unit copy, identify the cell and block
grids by their common center, straight cut ray and direction, origins,
and additive branch convention, and require the local cell and block
types on \(D_m(X)\) to match under that identification.
The construction below checks these
conditions summand by summand.  Its auxiliary \(P\)-boxes are
containers used for estimates and do not add inputs to an activity.

\begin{theorem}[Cutoff entry with one compulsory component]\label{ent:initialization}
Fix the finite list of map norms and reserves just described.  There
are constants \(C,c>0\) such that the following holds for every
sufficiently large fixed dyadic \(R\), with \(P=R^5\) and the fixed
period \(q\).  For every sufficiently large dyadic \(s\) satisfying
\eqref{ent:entry-side}, the assertions below hold uniformly in the
finite volume.  The threshold on \(s\) may depend on all preceding
choices.  A compatible torus or aligned free square has side at
least \(sP\); for a marked square its side is larger than the fixed
multiple of \(sP\) needed to contain the compulsory region and its
collars.  In the square applications its ratio to \(sP\) tends to
infinity.

\begin{enumerate}
\item There are plane activities \(K_m^0\) which are real on real
fields, negation-even, invariant under common translation by
\(\omega\), and covariant under block translations and square
symmetries.  There are centered torus and free-square activities
\(K_{\mathcal G,m}^0\) with
\[
                  \|K_m^0\|_{m,A_w,V},
                  \ \|K_{\mathcal G,m}^0\|_{m,A_w,V}
                         \le Ce^{-cR}
\]
in each of the fixed norms.  On every such finite graph a positive
field-independent scalar \(c_{\mathcal G,m}\) makes
\eqref{ent:ordinary-subset} hold.  If \(D_m(X)\) unwraps in a torus,
its activity is exactly the plane activity.  The same is true in a
free square when \(D_m(X)\) misses the wall.

\item In a free square with the center unit there are real periodic
activities \(K_{\mathcal G,m}^\tau\) and \(J_{\mathcal G,m}\) for which
\eqref{ent:labelled-subset} holds with that same scalar.  For a constant
\(C_R<\infty\) independent of \(s,n\),
\begin{equation}\label{ent:label-norm}
 \|K_{\mathcal G,m}^\tau\|_{m,A_w,V}\le Ce^{-cR},
 \qquad
 \|J_{\mathcal G,m}\|_{m,A_w,V}\le C_Rs^{C_R}.
\end{equation}
The sector norm may use the slightly weaker weight and stronger
regulator already reserved above.

There is a fixed set \(\mathfrak C_{\star,m}\) of at most \(C_R\)
entry blocks, containing the full inflation of a square of side a
fixed large multiple of \(P\) in observation-cell units around the
source, such that
\begin{equation}\label{ent:compulsory-support}
       J_{\mathcal G,m}(X)=0
                  \quad\hbox{unless}\quad
                  \mathfrak C_{\star,m}\subset X.
\end{equation}
Every block of \(\mathfrak C_{\star,m}\) is occupied in \(X\) in
the subset functional.  The full support of each summand of \(J(X)\)
also occupies every padded dependency and the complete connecting,
winding, or wall-reaching tag of every item in that summand.

\item If the rectangle spanned by \(D_m(X)\) misses the source, then
\begin{equation}\label{ent:translated-copy}
 K_{\mathcal G,m}^\tau(X,\psi)
             =K_{\mathcal G,m}^0(X,\psi+\lambda_X),
\end{equation}
where \(\lambda_X\) is its local lift.  This assertion uses the
finite-square centered prescription if the support is near a wall.
There are also centered-unit plane prescriptions \(K_{\infty,m}^\tau\)
and \(J_{\infty,m}\).  Whenever \(D_m(X)\) misses the square wall
under the centered aligned identification just specified, its
prescription is exactly the plane one with background translated by
\begin{equation}\label{ent:plane-background}
                         (v_n-v_\infty)/\alpha .
\end{equation}
This identity holds for both the unlabelled and labelled activities.
The sufficient regions in all three copy statements include the
padded supports and all occupied tags.  The plane centered-unit
prescriptions obey the same bounds as in \eqref{ent:label-norm},
with their local plane norms.
\end{enumerate}
\end{theorem}

The assertions concern arbitrarily late starts, with constants
uniform for all sufficiently large admissible \(s\).  They make no
uniform claim in temperature and assume no rate for a height
scaling limit.  In particular the smallness can be prescribed in
advance by increasing the fixed \(R\), while the label is only
required to obey \eqref{ent:label-norm}.

\subsection{The finite expansion and the occupied cover}

We prove Theorem~\ref{ent:initialization} by adapting the finite
expansion and pattern estimates in \cite[Section~9]{B}.
The theorem there is stated for quadratic discrete Gaussian heights;
we use its model-independent algebra, Gaussian estimates, and
counting, and prove the needed Bessel estimates with the cutoff.
The strengthened-law entry of \cite[Section~5]{C} is not an entry
theorem for the present law.  Its proofs of the bounded-data and
analytic estimates will be used with the inputs verified below.

In observation-cell units, start with the square grid of side \(R\).
Remove the square of side \(R/2\) centered at each vertex from its
incident squares and make each removed square a tile.  The trimmed
squares and these vertex tiles partition the cells.  At most three
tiles meet at a junction, and tiles meeting there share a face
segment.  At a free wall clip complete aligned arrays, terminate the
tiles at the wall, and use the same rules on every complete
neighborhood.  All pieces remain unions of whole cells and belong
to the rectilinear geometries of Proposition~\ref{prop:height-inputs}.
The four observation cells meeting the source lie with clearance
inside a vertex tile.  No tile deletion cuts through their unit
junction.

Use the finite colors and bounded proximity groups of
\cite[Section~9, \emph{A finite, covariant isolation expansion}]{B}.
For clarity, the colors specify the tile type and its square-symmetry
orbit modulo \(q\) in units of \(R\).  The period is greater than
the fixed physical proximity, core, open-proxy, guard, testing, and
skip radii of the isolation step.  Each proximity
group has a bounded number of members, and groups of one color have
disjoint enlarged physical neighborhoods used by the isolation
comparisons.  Order the colors.  At an event isolate its
group of tiles by deleting all still-retained bonds to other tiles,
and delete the corresponding cutoff faces.  Its core, open patch,
guard, and larger guard-testing patch have the fixed multiples of
\(R\) and separations in the cited construction.  The same coloring
and grouping apply to clipped wall groups.

Each open patch is a restriction of the current graph and therefore
keeps its earlier cuts and retained observation rows.  Measured by
center sup-distance from a group member in cell units, the cited
construction takes core tiles at distance at most \(R\), open-patch
tiles at distance at most \(5R\), and guard tiles in that patch
starting at distance at least \(3R\).  After the fixed stencil inflation of
\eqref{ent:guarded-face-error}, the core and changed artificial
boundary still have separation \(cR\), including for clipped or
disconnected groups.  Thus every remaining Gaussian
deletion-versus-patch difference estimated below is an instance of
that four-graph estimate.

If \(o_i,c_i\) are the current open patch before and after event
\(i\), set
\begin{equation}\label{ent:proxy}
 \mathcal B_i^\tau(y)=C_i
          \exp\{-\tfrac12(E_{o_i}^\tau(y)-E_{c_i}^\tau(y))\},
 \qquad C_i=\frac{Z_{o_i}^0(0)}{Z_{c_i}^0(0)}.
\end{equation}
Here \(C_i\) is the centered open-patch scalar; only the Gaussian
shape depends on the sector.  Replace the old
weight by \(\mathcal B_i^\tau\) times the cut weight plus their
difference.  Selecting the difference marks the event; after a
mark skip every later event within the fixed skip radius of its
group.  The radius is larger than the physical core, open-proxy,
guard, and testing collars of the isolation rules.  A later considered patch then misses each earlier marked
core and guard.  Its proxy does not depend on the choice of either
signed term inside that earlier difference.  This is a finite
binary identity at the level of the integrand.  After the last
color, all undeleted cross-tile bonds lie in fixed halos of marks.

Let the superscript \(0{\rm m}\) denote the path with no marks.
Dividing the expansion by
\[
 \prod_i\mathcal B_i^{\tau,0{\rm m}}(y)
              \prod_b Z_b^0(0)e^{-E_b^\tau(y)/2}
   =\left(\prod_i C_i^{0{\rm m}}\prod_b Z_b^0(0)\right)
                   e^{-E_{\rm asm}^\tau(y)/2}
\]
defines the assembled Gaussian energy \(E_{\rm asm}^\tau\).
The index \(b\) runs over the isolated tiles on that path.
The scalar on the right is positive and is exactly the same in
the centered and unit calculations.  At a real zero centered
observation every retained cutoff equals one, so these scalar
normalizers need no cutoff correction.  In a marked halo retain
the whole signed height factor and every affected proxy.  Outside
the halos expand each isolated factor as one plus
\[
 T_Q^\tau(y)-1,\qquad
 T_Q^\tau(y)=
          \frac{\widehat Z_Q^\tau(y)e^{E_Q^\tau(y)/2}}
                    {Z_Q^0(0)} .
\]
Selecting this second term selects an additional tile.  In every
term of the unit calculation insert one formal compulsory marker
\(I_\star\) at the source.  If the source tile is isolated, leave
its normalized tile factor unexpanded and attach it to this marker.
If it lies in a marked halo, the marker carries the factor one and
that halo retains its entire signed factor and its ordinary marked
items.  The marker is adjacent to them and hence labels their
component.  In particular, an arbitrarily long chain of attached
halos keeps the union of their individual covers; it is not hidden
inside the fixed primitive cover of \(I_\star\).  Every term has
exactly one marker.  Identical proxies and final tile factors off
the affected halos cancel exactly, before any estimates are made.

The remaining Gaussian mismatch is
\((E_{\mathcal G}^\tau-E_{\rm asm}^\tau)/2\).
Telescope it over the deletions in the no-mark path.
The \(D_s\sigma_z^2\) terms of \eqref{ent:face-decomposition}
telescope exactly, since a full junction has the same \(D_s\)
in all graphs in which it persists.  Lemma~\ref{ent:face-localization}
therefore expresses the mismatch as a quadratic sum in face
values whose coefficients, after summing deletion positions, obey
\[
                         C e^{-cR-c\,d(f,f')}.
\]
The retained core-distance decay in
\eqref{ent:guarded-face-error} justifies this summation with a
constant independent of volume.  Enumerate its diagonal and
off-diagonal summands as \(b_\ell w_fw_{f'}\), with all symmetry
factors included in \(b_\ell\), and expand
\begin{equation}\label{ent:face-links}
 e^{(E_{\mathcal G}^\tau-E_{\rm asm}^\tau)/2}
       =\prod_\ell(1+\ell_\ell),
       \qquad \ell_\ell=e^{b_\ell w_fw_{f'}}-1.
\end{equation}
If \(d_\ell\) is the full tagged length of the link, then
\(|b_\ell|\le Ce^{-c(R+d_\ell)}\).  The elementary inequality
\(|e^z-1|\le |z|e^{|z|}\), followed by
\(2|w_fw_{f'}|\le w_f^2+w_{f'}^2\), gives a factor
\(Ce^{-c'(R+d_\ell)}\) times a quadratic exponential with
coefficients \(Ce^{-c'(R+d_\ell)}\) on its incident faces.
For example choose those coefficients proportional to
\(|b_\ell|^{1/2}\), which pays the polynomial prefactor while
retaining half the exponential gain.  Such costs may also be
charged on its connecting path.  The sum of all reserved
coefficients at one face is \(Ce^{-c'R}\), since there are only
polynomially many endpoints and path descriptions at a given
length.  At a fixed finite real observation the product in
\eqref{ent:face-links} is absolutely convergent; the integrated
absolute convergence will follow from the estimates below.

We retain the following geometric information in a link.  On a
torus, periodize the columns of the massive observation inverse
and retain each displacement image separately.  This is an exact
periodization because the columns have the exponential decay
proved above and the local forcing and proxy rules are periodic.
The link's path keeps its entire winding.  In a free square use
the plane coefficients between retained faces, with coordinate
paths inside the square, and put their difference from the square
matrix into additional links.  Give each additional link a path
that also reaches a nearest wall.  The cut-distance bound in
\eqref{ent:face-matrix-bounds}, together with the direct
face-distance and \(R\)-gains, bounds this coefficient by
\(Ce^{-c(R+d_\ell)}\) with the wall path included in \(d_\ell\).
The fixed \(O(R)\) proxy ranges are absorbed by decreasing \(c\).
Average coordinate-order choices and split ties to preserve
symmetry.  Each path has linear length, so its cover uses
\(O(1+d_\ell/R)\) entry blocks; no bounding rectangle is filled.
This decomposition makes interior copying exact, since every
term that differs at a wall has a tag reaching that wall.

We now specify the occupied supports in the subset functional.
For an ordinary mark or selected tile \(I\), first choose one fixed
physical envelope containing its tile
halo and every core, open-proxy, guard, testing, and seam-comparison
piece used for any allowed translate.  The envelope has \(O(R^2)\)
observation cells per item, with a constant that may depend on the
fixed \(q\); every later seam comparison is
clipped to this envelope.  Some cells in the envelope serve only the
estimates; occupying them is a fixed conservative choice made before
the later box translate is optimized.  Add all fixed field and Gaussian-averaging
collars of the local calculus.  For a link
include its entire chosen connecting path and every winding or
wall-reaching tag, with the same collars.  All these sets are
enlarged to whole observation cells.  Let
\(\operatorname{cov}_m(I)\) be all \(m\)-blocks met by this full
set.  The collars include the finite range of \(C_{<m}\) where
needed for the averaging dependency.  They are fixed once the
map parameters and \(q\) are fixed.  A halo cover has bounded
entry-block size, and a link cover has size
\(O(1+d_\ell/R)\).
The Gaussian-averaging collars can overlap even for events of one
color; the full-cover adjacency below merges every such dependency.

For the compulsory source item, let the compulsory region be the
observation-cell square within a fixed large multiple of \(P\) of the
center.  Begin with this region and the local source-tile dependency
when that tile factor is attached.
Apply the same full inflation, including every padded dependency
of that local factor and any tag assigned to it, and call the
resulting set of entry blocks
\(\mathfrak C_{\star,m}\).  Choose the multiplier of \(P\) large
enough that, for any item outside the compatibility neighborhood of
this cover, every allowed \(P\)-box containing part of its physical
region, together with its fixed testing enlargement, misses the source
and uses no stencil near it.  This is possible because such a box and
enlargement reach only \(O(P)\) cells from the item, whereas the initial
square may use a larger fixed multiple of \(P\).  Its size is at
most \(C_R\), because \(m=qRs\) and \(P,R,q\) are fixed as
\(s\) increases.  It is connected and contains the center.

Declare two items adjacent whenever their full covers are within
the fixed connection distance of the subset calculus.  Its initial
choice includes the Gaussian range and the binary expansion's
exclusion rules: \(C_{<m}\) has range less than \(32m\) for the
printed polynomial, and the exclusions have diameter
\(O(Rs)=O(m/q)\).  Thus one bound chosen before \(L\) works for
every later \(R,s\).  This distance is fixed
independently of the size of the compulsory cover; the large
cover occupies more blocks and does not enlarge the connection
radius used by subsequent maps.  For a component \(\pi\) without the
compulsory item, define its occupied support to be the full union
\begin{equation}\label{ent:occupied-cover}
             X(\pi)=\bigcup_{I\in\pi}\operatorname{cov}_m(I).
\end{equation}
For a component \(\pi_\star\) containing the compulsory item, define
its occupied support to be exactly
\begin{equation}\label{ent:occupied-label-cover}
 X(\pi_\star)=\mathfrak C_{\star,m}
       \ \cup\!
       \bigcup_{\substack{I\in\pi_\star\\I\ne I_\star}}
                        \operatorname{cov}_m(I).
\end{equation}
In particular, a summand of \(J(X)\) is assigned only to
\(X=X(\pi_\star)\).  The full set \(\mathfrak C_{\star,m}\) is
occupied even where the source factor is constant in the field.
Every padded dependency and every tag of an attached item is
occupied as well.  We never replace this set by a root block or
by a smaller subcover.  Items that meet any part of this
compatibility neighborhood belong to \(\pi_\star\), even if
their factors would otherwise be independent.  Thus their
ordinary patterns have not been discarded; they are summed in
the labelled component.

Separated components now have separated field and noise
coordinates and separated Gaussian subvectors below \(m\).
They also have no exclusion constraint between them: every
such constraint was included in the adjacency just defined.
Within a component, skip and validity constraints are functions of
its recorded items and local rule data.  Potentially incompatible
exterior items need not lie in \(D_m(X)\); when selected, their
exclusion is enforced by subset compatibility rather than consulted
as an input to this component's activity.
Averaging a pattern in \eqref{ent:averaged-interaction} therefore
factors over these components.  The sum over components without
the compulsory item defines the ordinary activity; the sum over
the component \(\pi_\star\) defines \(J\).  All constraints within
a component remain in its defining sum.  This gives the subset
algebra \eqref{ent:ordinary-subset} and
\eqref{ent:labelled-subset}, once the absolute convergence below
justifies the exchanges with Gaussian integration.  The common
scalar is the positive scalar of the no-mark divisor divided by
\(Z_{\mathcal G}^0(0)\).  No independence of the random guards
or of the binary marks has been assumed.

In the unit calculation, define \(K^\tau(X)\) by its ordinary
pattern sum whenever \(X\) can coexist with
\(\mathfrak C_{\star,m}\).  Values on supports that cannot
coexist do not enter \([x]\cZ_m(K^\tau+xJ)\), by
\eqref{ent:occupied-label-cover}.  Set those values to zero,
except that when the rectangle spanned by \(D_m(X)\) misses the
source, set them to the translated centered prescription.
This convention is useful for later local maps and does not
alter the exact identity.

We check the sufficient copy region at the level of each summand
assigned to \(X\).  Every evaluated field or noise coordinate and
every local isolation or proxy stencil lies in the full cover and
hence in \(D_m(X)\).  The grid phase and color rule is common under
the compatible aligned identification.  The cover also contains the
whole averaging collar for \(C_{<m}\), whose range is less than
\(32m\).  It occupies every projected path site, while the summand
retains the complete lifted winding and tag data.  Compatible aligned block
identifications preserve these local data and the within-component
constraints.  The optimized \(P\)-grid, its energy classification,
and empty portions of its boxes are used only in the proof of the
bound; they are not data in the defining activity.

There is one qualification for the numerical link coefficients.
An untagged coefficient is the fixed universal plane table indexed
by its faces and displacement; the definition of that table may use
the plane inverse.  A torus image carries its complete winding path.
A square-minus-plane coefficient may depend on the whole square,
but its additional link carries a complete path to the wall.
Consequently a summand with such an ambient obstruction cannot be
assigned to an \(X\) whose \(D_m(X)\) unwraps, or respectively misses
the wall.  All contributing rules and coefficients in that case are
the common plane ones.  This proves the centered copy statements.

In a source-free rectangle spanning \(D_m(X)\), an integral gauge
removes \(\tau\).  It turns every retained factor in
\eqref{ent:cutoff-definition} into the ordinary neighbor cutoff and
changes the field argument to the whole lifted translate
\(\psi+\lambda_X\), proving \eqref{ent:translated-copy}.  Near a wall
the two sectors use the same finite-square centered coefficient,
including its possible dependence on remote parts of that square.
Under the centered aligned identification in the theorem, when
\(D_m(X)\) misses the wall, the affine energies, cutoff, and source
factor are the same plane-unit rules.  Changing
\(v_\infty\) to \(v_n\) translates their argument by
\eqref{ent:plane-background}.  The full compulsory cover and every
attached cover remain occupied under these identifications.  The
artificially filled entries above are either zero or the same
translated centered prescription.  Thus the identities hold for
the complete component sums with \(D_m(X)\) as a sufficient entry
copy region.  The later causal neighborhood in
Section~\ref{sec:observable} tracks the subsequent maps separately.

\subsection{Bounds for the pattern sums}

Fix a pattern \(\pi\).  Let \(j_1\) be its number of binary marks,
\(j_2\) its number of additional selected tiles, and \(d_\ell\)
the tagged lengths of its links.  Let \(X\) be the full occupied
cover just defined.  Write \(F_\pi(\psi)\) for its signed factor
after averaging the noise and the Gaussian pieces below \(m\),
with the common scalar removed.  Put
\[
\mathcal C_\pi(s)=
 \begin{cases}
  1,&\pi\text{ has no compulsory item},\\
  C_Rs^{C_R},&\pi\text{ has the compulsory item}.
 \end{cases}
\]
The constants \(C_R\) in this notation may increase during the proof,
but remain independent of \(s\) and the volume.

\begin{lemma}[The analytic pattern estimate]\label{ent:pattern-estimate}
After the parameter choices in \eqref{ent:reserve-order} and
\eqref{ent:entry-side}, and the threshold choice described below,
there is \(s_0<\infty\) such that for every admissible \(s\ge s_0\),
\begin{equation}\label{ent:analytic-pattern-bound}
 |F_\pi|_{m,X,\psi}
 \le \mathcal C_\pi(s)\,
       C^{j_1+j_2}e^{-cR(j_1+j_2)}
       \prod_\ell Ce^{-c(R+d_\ell)}
       W_m^\kappa(X,\psi)e^{V_{j_0}(X,\psi)}
\end{equation}
for all real backgrounds \(\psi\).  The point norm on the left is
the full derivative sum in \eqref{eq:point-norm}.  The constants
are uniform in the compatible tori and free squares.
\end{lemma}

\begin{proof}
We first obtain a real-field bound, keeping the physical comparison,
the finite conditional limits, and the Gaussian averaging as separate
steps.

\paragraph{Centered normalization on the affected halos.}
Prescribe the primary guard data \(b_i^p\) at every marked event.
Let \(K_o^\tau,K_c^\tau\) be its two positive interior kernels,
including their retained cutoff factors.  For real observations set
\[
 \delta_i(y,b_i^p)=
 \frac{|K_o^\tau(y,b_i^p)-\mathcal B_i^\tau(y)K_c^\tau(y,b_i^p)|}
      {K_o^\tau(y,b_i^p)+\mathcal B_i^\tau(y)K_c^\tau(y,b_i^p)}
                         \le1.
\]
Factors outside the guard cancel in this ratio.  Conditional
integration of a product of differences and the triangle inequality
give the sum of its positive topologies, with \(\prod_i\delta_i\)
under their guard laws.  This is valid with correlated guard data.

For a positive choice \(G\), let
\(\mathsf B_{\pi,G}^\tau(y)\) denote its nonnegative normalized halo
factor before the noise and Gaussian entry averaging.  It retains
the affected proxies and the conditional \(\delta_i\) product, and
uses the quotient left after cancelling the identical no-mark proxy
and tile factors off the affected halo union.  The intact selected
tiles, an intact isolated source tile when present, and the links
remain separate.  The triangle inequality bounds the normalized
unaveraged pattern by the finite sum of these halo factors times
the absolute values of those remaining factors; averaging this
bound gives a bound for \(F_\pi(\psi)\).  If there is no halo, the
single empty halo factor is one.  We now estimate one
\(\mathsf B_{\pi,G}^\tau\).

Only a number of centered normalizations proportional to the number
of affected halos is charged.  Here is the finite accounting.  In
the union of the halos take a further collar beyond every skip and
proxy radius, and give its artificial boundary a free wall.  For
either the positive branch or the no-mark branch, let \(T_e\) be
the true centered before-to-after ratio at a deletion in this
union:
\[
 T_e=\frac{Z_{\mathrm{before},e}^0(0)}
           {Z_{\mathrm{after},e}^0(0)}.
\]
These are zero-observation sums with every retained observation
row and no temporary hard guard.  Let \(\widetilde C_e\) be the
centered zero-observation proxy clipped to the union.  Exact
telescoping gives
\[
 Z_{\rm final}^0(0)\prod_e\widetilde C_e
           =Z_{\rm initial}^0(0)\prod_e\frac{\widetilde C_e}{T_e}.
\]
The initial zero-observation sum is the same in the two histories.  They agree
near the artificial boundary, so the clipped proxies there and
all final components off the affected union cancel in the
quotient.  There are \(O(j_1)\) remaining comparisons, with a
fixed additional number depending on \(R\) if the compulsory
item is present.  Each centered restoration is between its open
patch value and its value with a separating zero guard.  The
logarithm of this gap is bounded by \(D(H,p)\): a full primary
pin on \(H\) containing the restored endpoints makes the
restoration a fixed scalar.  Proposition~\ref{prop:comparisons}
gives the direction, and Proposition~\ref{prop:height-inputs}
and the Gaussian localization in \cite[Section~3]{B} bound the
gap on each of these finitely many shapes.  This is the local
double telescoping of \cite[Section~9, Lemma~9.3]{B} and
\cite[Equation~(4.6)]{C}, applied to the present centered
normalizers.  No comparison is made over the unaffected volume.

Apply the identical accounting to the Gaussian shapes.  The
bounded face matrices in Lemma~\ref{ent:face-localization}
give \(Ce^{-cR}\mathcal E_\pi(w)\), where
\(\mathcal E_\pi\) is the sum of \(w_f^2\) on the underlying
faces in the original padded halos.  The singular terms cancel
whenever their junction persists.

For each resulting positive halo topology \(G\), we now name the
positive factor and its guard law.  If \(p\) is the current set of
primary guards and \(\gamma\) their values, let
\begin{equation}\label{ent:joint-density}
\begin{split}
 \mathcal J_G^\tau(y,\gamma)
  &=\sum_{h_p=\gamma}\mathcal W_G^\tau(h)
               \exp\{-\tfrac12\|B_sh-y\|^2\},\\
 Z_G^\tau(y)&=\sum_\gamma\mathcal J_G^\tau(y,\gamma),\qquad
 \nu_{G,\tau,y}^p(\gamma)
       =\frac{\mathcal J_G^\tau(y,\gamma)}{Z_G^\tau(y)}.
\end{split}
\end{equation}
Here \(\mathcal W_G^\tau\) is the unnormalized product of the
retained primary height weights.  The common Gaussian density
constants are omitted.  Thus \(\mathcal J\) is an unhatted joint
observation and guard density, distinct from the labelled activity
\(J\).  Write \(\E_{G,\tau,y}\) for expectation under this guard law.
With
\begin{equation}\label{ent:positive-factor}
\begin{split}
 \Xi_G^\tau(y)&=\prod_{f\in\mathcal F(G)}\chi_s(w_f),\\
 \mathcal A_G^\tau(y)
 &=e^{E_G^\tau(y)/2}\frac{\widehat Z_G^\tau(y)}{Z_G^0(0)}
 =e^{E_G^\tau(y)/2}\Xi_G^\tau(y)\frac{Z_G^\tau(y)}{Z_G^0(0)},
\end{split}
\end{equation}
the cutoff \(\Xi_G^\tau(y)\) depends only on the prescribed
observation and does not alter the guard law.  The centered and
Gaussian accounting just proved gives
\begin{equation}\label{ent:positive-topology}
\begin{split}
 \mathsf B_{\pi,G}^\tau(y)\le {}&
 e^{Cj_1+C_R1_{\{\star\}}+Ce^{-cR}\mathcal E_\pi(w)}\\
 &\quad{}\times\mathcal A_G^\tau(y)
          \E_{G,\tau,y}\prod_i\delta_i .
\end{split}
\end{equation}
The factor \(\mathcal A_G^\tau\) is retained in every halo topology.
In particular it contains any surviving unit self energy there.
The notation \(1_{\{\star\}}\) means that the pattern has the
compulsory item.

Before making the box cuts, let \(D_{\mathrm{phys}}(\pi)\) consist of
the observation cells and original pre-deletion faces used by the
pattern's positive halo and selected-tile expressions, their affected
open proxy patches, and the fixed core, guard, testing, and seam collars
used in these comparisons, all inside the per-item physical envelopes
occupied above.  Include the original selected-tile region
also when that expression takes its constant-one branch, and include
the local source-tile region when it is attached.  Each ordinary mark
or selected tile contributes \(O(R^2)\) cells with these fixed collars;
the source region contributes only a cost depending on \(R\).
The auxiliary boxes used below are containers only and add no empty
cells or faces to \(D_{\mathrm{phys}}(\pi)\).
The complete link paths and the collars used only for their covariance
dependence remain in the full occupied cover \(X\) and in \(D_m(X)\).
Artificial occupancy of the compulsory cover with marker factor one
creates no additional physical \(Z\)-factor.

\paragraph{A real cut of the halos into \(P\)-boxes.}
Cut only the positive halo topology \(G\) into boxes of
side \(P=R^5\) in cell units, on a translate of the \(R\)-grid.
This grid clips its physical region inside the fixed envelopes.
Each resulting graph contains
only its actually present cells and faces there; the cut never
introduces field dependence on a remote part of a \(P\)-box.
On a torus remove every bond crossing a lifted box seam, including a
bond between two lifted copies of the same projected box when the
torus side is exactly \(sP\).
Retain a mark for its conditional gain only if its guard and
larger guard-testing patch lie deep inside a box.  Averaging the
finitely many translates loses at most \(C(R/P)j_1\) marks.
For the bounded face matrices, the corresponding average
seam cost is at most \(C(R/P)\mathcal E_\pi(w)\).  To see this,
let \(\Sigma\) be the altered seam with its fixed comparison collar.
Combine the two decays in \eqref{ent:face-matrix-bounds}.
The inequality \(d(f,\Sigma)\le d(f,f')+d(f',\Sigma)\),
row summation, and \(2|w_fw_{f'}|\le w_f^2+w_{f'}^2\) bound
the cost by
\[
 C\sum_{f\text{ in the original padded halos}}
                  e^{-c\,d(f,\Sigma)}w_f^2.
\]
The remaining pair-distance decay makes the row sum finite.
Averaging the \(R\)-spaced translates of the \(P\)-grid gives
\(\operatorname{avg}_{\Sigma}e^{-c\,d(f,\Sigma)}\le CR/P\)
at every such face.  Only these existing halo faces are charged.
The simultaneous Markov bound for the lost marks
and this nonnegative quadratic cost selects a translate
retaining at least half of the marks, for large \(R\), with
the stated cost multiplied by a fixed constant.  If a seam
removes the unit junction, \eqref{ent:fixed-circulation}
adds \(C\log(2+s)\) once at the compulsory component.

For the physical comparison, first discard the factors
\(\delta_i\le1\) at lost guards and let \(p\) denote the retained
deep guards.  Write \(G_{\rm cut}=\bigsqcup_bG_b\) for the disjoint
union of the actual clipped halo pieces.  At every real \(y\) and
every retained guard value \(\gamma\), the affine-to-centered
comparison gives the joint-density inequality
\begin{equation}\label{ent:seam-joint}
 \frac{\mathcal J_G^\tau(y,\gamma)}{\mathcal J_G^0(0,0)}
 \le
 \frac{\mathcal J_{G_{\rm cut}}^\tau(y,\gamma)}
      {\mathcal J_{G_{\rm cut}}^0(0,0)}.
\end{equation}
Indeed restoring the seam bonds adds the same positive matching
precision in the affine and centered height, noise, and guard
coordinates.  Proposition~\ref{prop:comparisons} applies to this
fiber also when the sector displacement is nonzero.  The joint
densities in \eqref{ent:seam-joint} contain no cutoff factors.

Put
\[
 \mathsf p_G^0=\frac{\mathcal J_G^0(0,0)}{Z_G^0(0)},\qquad
 \mathsf p_{\rm cut}^0=
       \frac{\mathcal J_{G_{\rm cut}}^0(0,0)}{Z_{G_{\rm cut}}^0(0)}.
\]
These are the centered zero-guard probabilities at zero observation.
For each actual halo piece define the complete per-box factor
\begin{equation}\label{ent:box-factor}
 Q_b^\tau(y_b)=\mathcal A_{G_b}^\tau(y_b)\,
       \E_{G_b,\tau,y_b}\prod_{i\ {\rm deep\ in}\ b}\delta_i .
\end{equation}
The expectation is one when there is no retained guard in \(b\).
The energies, cutoff products, centered normalizers, and joint guard
sums on \(G_{\rm cut}\) factor over its actual pieces.  Sum
\eqref{ent:seam-joint} against the unchanged nonnegative product of
the retained \(\delta_i\), and multiply by
\(e^{E_G^\tau(y)/2}\Xi_G^\tau(y)\mathsf p_G^0\).  This gives
\begin{equation}\label{ent:seam-normalized}
\begin{split}
 \mathcal A_G^\tau(y)\E_{G,\tau,y}\prod_{i\ {\rm retained}}\delta_i
 \le {}&
 e^{(E_G^\tau(y)-E_{G_{\rm cut}}^\tau(y))/2}
 \frac{\Xi_G^\tau(y)}{\Xi_{G_{\rm cut}}^\tau(y)}\\
 &\quad{}\times
 \frac{\mathsf p_G^0}{\mathsf p_{\rm cut}^0}
 \prod_b Q_b^\tau(y_b).
\end{split}
\end{equation}
The local kernel at a retained guard is unchanged because its
testing patch lies inside the box.  The cutoff quotient in
\eqref{ent:seam-normalized} is at most one on real data.  Add the
zero guards one at a time and bracket each by its common open and
outer-zero-guard values.  The separated testing patches and the
same pin-interaction bound used above give
\[
       \frac{\mathsf p_G^0}{\mathsf p_{\rm cut}^0}
          \le e^{Cj_1+C_R1_{\{\star\}}}.
\]
The optimized shape difference is bounded by the previously
calculated \(C(R/P)\mathcal E_\pi(w)\), with
\(C1_{\{\star\}}\log(2+s)\) if the seam removes the unit junction.
The earlier proxy cost \(Ce^{-cR}\mathcal E_\pi(w)\) in
\eqref{ent:positive-topology} remains separate.  This is the
sector seam estimate of \cite[Equation~(9.21)]{B}, now with its
normalization and per-box factors displayed.

For each real \(y\), choose the first qualifying translate in a
fixed order of the finite allowed list.  In the sum over that list,
keep the indicator of this \(y\)-dependent choice until the bounds
for retained marks and seam costs have been applied.  Then discard the
indicator, which is at most one, and perform the Gaussian integral
for each translate.  This is a sum over only polynomially many
halo cut choices.  A nonempty ordinary halo pattern pays that
polynomial from its retained mark gains, and a labelled pattern may
also pay it in \(C_R\).

Keep every isolated selected tile \(Q\) intact and assign it to its
own auxiliary box containing \(Q\), of side at most \(R\le P\).
These boxes are independent of the translated halo grid; the entire
tile, including faces crossed by a halo seam, stays in one group.
This is the separate isolated-tile grouping of \cite[Section~9]{B}.
The actual isolated tile graphs are mutually disjoint and lie outside
the actual halo graph \(G\), even when their comparison or dependency
collars overlap.  If the source tile is isolated and unexpanded, keep
it as one intact compulsory group as well.  When the source lies in a
halo, its marker has factor one and adds no second physical group.
Patterns with only isolated tiles, links, or an isolated source marker
need no halo cut choice.

\paragraph{Box energies and their classification.}
Let \(b\) range over the nonempty clipped halo boxes and the auxiliary
tile boxes.  For a halo box, let \(\mathcal E_b(w)\) be the sum of
\(w_f^2\) on the original pre-deletion faces of its physical halo
region after clipping to \(b\).  For the auxiliary box of an isolated
tile \(Q\), put
\[
                  \mathcal E_b(w)=\sum_{f\in\mathcal F(Q)}w_f^2,
\]
using every original face of the intact tile.  Define this energy
before bounding \(T_Q^\tau-1\) by its positive and constant-one
branches, so both branches have the same classification.  The current
halo topology can have fewer faces than its underlying region.
Neither energy uses empty container cells.  A box met solely by a
link cover is charged in the separate link bounds.  A box is good when
\(\mathcal E_b(w)\le U_{\rm big}^2\), where \(U_{\rm big}\)
is a finite threshold still to be chosen.  In a pattern containing
\(I_\star\), call a box \emph{compulsory} if its fixed testing
enlargement meets the compulsory region; the intact isolated source
tile is compulsory as well.  There are only a number depending on
\(R\) of these boxes and of their possible marks and selected tiles.
Every other box uses the ordinary good or bad rule.  In particular,
an unlabelled unit-sector pattern uses those rules even if a container
overlaps the compulsory region.  Its full cover can coexist with
\(\mathfrak C_{\star,m}\), so the size choice above makes each such
box and its testing enlargement source-free.  Thus every ordinary
good box admits the local gauge used next.

\paragraph{Ordinary good boxes and their conditional gain.}
In an ordinary good box use an integral gauge on its source-free
rectangle.  Even a pattern that winds globally admits this local
lift.  Translate each component of the underlying pre-deletion graph
by a common height period so that one observed value is in
\([-\pi,\pi]\).  Paths of length bounded by a polynomial in \(P\),
together with the good-box bound, put all lifted observations in a
fixed compact set depending on \(R\) and \(U_{\rm big}\).  There are
finitely many box geometries, topologies, and mark lists at these
fixed parameters.  Proposition~\ref{prop:height-inputs}, after
convolution with the unit observation noises, gives uniformly on
these compact sets
\[
                    \mathcal A_{G_b}^\tau(y_b)\longrightarrow1 .
\]
Here \(G_b\) is an actual clipped halo piece or an intact tile graph,
with the local gauge understood.  The finitely many cutoff logarithms
tend uniformly to zero on the compact face data by
\eqref{ent:cutoff-tails}.  Thus the positive normalized topology
alone is bounded by a fixed constant on every ordinary good box for
all sufficiently large \(s\).  We next estimate the conditional
factor in \(Q_b^\tau\) for a halo box.

Fix one retained mark \(i\), and write \(p=p_i\) for its full primary
guard.  Let \(t=o,c\) be the two complete adjacent box topologies
differing only at this core, with every other positive choice fixed.
Use the same component-period recentering in the two topologies and
write \(y^\ell,\gamma\) for the lifted observation and guard data.
Both topologies keep the same observation rows.  In this paragraph
\(\mathcal J_t(y^\ell,\gamma)\) is the joint density
\eqref{ent:joint-density} in the centered gauge, and
\(\E_{t,0}\) samples the centered full-box law at zero observation.
Set
\begin{equation}\label{ent:one-mark-law}
\begin{split}
 z_t(\gamma)&=\mathcal J_t(0,\gamma),\qquad
 Z_t(0)=\sum_\gamma z_t(\gamma),\\
 L_t(y^\ell;\gamma)
   &=\E_{t,0}\!\left[e^{\langle y^\ell,B_sh\rangle}
                                      \mid h_p=\gamma\right],\\
 \mathcal J_t(y^\ell,\gamma)
   &=e^{-\|y^\ell\|^2/2}z_t(\gamma)L_t(y^\ell;\gamma),\\
 \nu_{t,y^\ell}^p(\gamma)
   &=\frac{z_t(\gamma)L_t(y^\ell;\gamma)}
           {\sum_g z_t(g)L_t(y^\ell;g)}.
\end{split}
\end{equation}
The last line is the sampled guard marginal of the centered full-box
law tilted by \(e^{\langle y^\ell,B_sh\rangle}\).  The cutoff is a
function of the observation only and does not change this law.
The box normalizers \(Z_t(0)\) in \eqref{ent:one-mark-law} are
distinct from the open-patch scalar
\(C_i=Z_{o_i}^0(0)/Z_{c_i}^0(0)\) in \eqref{ent:proxy}.

Define the centered conditional restoration ratio
\[
 r_\Delta(\gamma)=
 \frac{z_o(\gamma)/z_c(\gamma)}{z_o(0)/z_c(0)}.
\]
Factors exterior to the full guard cancel from the kernel quotient.
The exact one-mark identity is
\begin{equation}\label{ent:one-mark-identity}
\begin{split}
 \log\frac{K_o^\tau(y,b_i^p)}
              {\mathcal B_i^\tau(y)K_c^\tau(y,b_i^p)}
 &=c_{i,s}+\log r_\Delta(\gamma)+\Lambda_{i,s}(y,\gamma)\\
 &\quad{}+\log\frac{\Xi_o^\tau(y)}{\Xi_c^\tau(y)},\\
 c_{i,s}&=\log\frac{z_o(0)}{C_i z_c(0)},\\
 \Lambda_{i,s}(y,\gamma)
 &=\log L_o(y^\ell;\gamma)-\log L_c(y^\ell;\gamma)\\
 &\quad{}+\tfrac12\bigl(E_{o_i}^\tau(y)-E_{c_i}^\tau(y)\bigr).
\end{split}
\end{equation}
Here \(\Xi_o^\tau/\Xi_c^\tau\) is the quotient of the retained local
cutoffs in the two kernels.  The energy term in \(\Lambda_{i,s}\)
is \(-\log(\mathcal B_i^\tau/C_i)\).  The first term \(c_{i,s}\)
compares the centered zero-observation, zero-guard restoration with
the open-patch proxy.  The separated pin interaction and its
Gaussian localization give a nonnegative upper bound
\[
 |c_{i,s}|\le d_{i,s}:=D_i(H,p),\qquad
 d_{i,s}\longrightarrow d_{i,{\rm g}}\le Ce^{-cR}.
\]
Here \(D_i(H,p)\) is the centered pin interaction for this
separated comparison.  This bounds the restoration gap without
asserting a limit for \(c_{i,s}\).
The local cutoff logarithm is uniformly \(o(1)\) on compact lifted
data.

The \(r_H\) of Proposition~\ref{prop:height-inputs}, formed in the
centered cut law with a full pin on a core containing the restored
endpoints and the supports of all restored forms, satisfies
\(r_H\le r_\Delta\le1\).  These are centered
objects even when the guard is sampled under \(\nu_{t,y^\ell}^p\).
Indeed the exact change of the sampled guard law is
\begin{equation}\label{ent:tilted-guard-change}
 \frac{\dd\nu_{o,y^\ell}^p}{\dd\nu_{c,y^\ell}^p}(\gamma)
 =
 \frac{r_\Delta(\gamma)L_o(y^\ell;\gamma)/L_c(y^\ell;\gamma)}
 {\E_{\nu_{c,y^\ell}^p}
       [r_\Delta L_o(y^\ell;\cdot)/L_c(y^\ell;\cdot)]}.
\end{equation}
At zero observation this reduces to
\(r_\Delta/\E_{\nu_{c,0}^p}r_\Delta\).  For a nonzero observation
the Laplace quotient in \eqref{ent:tilted-guard-change} is retained.

Here is the quantitative Gaussian estimate for the terms in
\eqref{ent:one-mark-identity}.  Let \(u_c\) be the cut Gaussian
harmonic prediction from the guard, and let \(\chi_H\) equal one at
the core and vanish before the guard.  In the gradient-plus-unit-
observation energy norm \(T\), put \(Y_i=\|\chi_Hu_c\|_T\).  The
localized prediction estimate in \cite[Equation~(9.23)]{B} gives,
for every \(r\ge2\),
\begin{equation}\label{ent:prediction}
 \|Y_i\|_{L^r}
       \le C\sqrt r\,\operatorname{poly}(P)e^{-cR}
                         (1+\sqrt{\mathcal E_b(w)}).
\end{equation}
It holds under either topology's Gaussian guard law: restoration
decreases the centered guard covariance, and the fixed tilt adds a
deterministic prediction with the same localized operator bound.
The negative logarithm of the Gaussian full-core pin ratio is at
most \(Y_i^2/2\).  Subtracting this cutoff prediction also makes
the cut minimizer admissible for restoration without changing its
guard.  Orthogonal projection bounds the linear conditional-Laplace
discrepancy by \(C\|y^\ell\|_2Y_i\); the conditional quadratic
discrepancy from the open-patch proxy is
\(Ce^{-cR}\|y^\ell\|_2^2\).  These are finite-energy prediction
estimates and require no pointwise continuum trace.

We spell out the fixed-order passage to the sampled primary guards.
The elementary inequality
\[
 \left|\frac{x-1}{x+1}\right|\le\min(1,|\log x|),\qquad x>0,
\]
and \(r_H\le r_\Delta\le1\) give the pointwise bound
\[
\begin{split}
 \delta_i&\le\Theta_{i,s},\\
 \Theta_{i,s}&:=
 \min\left(1,\ d_{i,s}-\log r_H+|\Lambda_{i,s}|
                    +\left|\log\frac{\Xi_o^\tau}{\Xi_c^\tau}\right|\right).
\end{split}
\]
Fix \(R\), \(P=R^5\), a finite \(U_{\rm big}\), the finite family of
box geometries, topologies, and marks, a finite integer \(r_{\max}\) at least
two and at least every deep-mark count in that family, and
\(\eta>0\).  There is a
threshold \(s_0\), depending on these fixed quantities, such that
for all admissible \(s\ge s_0\), all compact lifted good observations,
either \(t=o,c\) sampled under \(\nu_{t,y^\ell}^p\), and every integer
\(2\le r\le r_{\max}\),
\begin{equation}\label{ent:one-mark-moment}
 \|\delta_i\|_{L^r(\nu_{t,y^\ell}^p)}
 \le C P^C r^C e^{-c_0R}(1+\mathcal E_b(w))^C+\eta .
\end{equation}
The constants and degrees in the Gaussian term are independent of
\(U_{\rm big},s,r\); the threshold \(s_0\) may depend on every
preceding fixed choice.

To prove this assertion, first fix the smooth guard lists and smooth
core penalties in Proposition~\ref{prop:height-inputs}, then pass to
the lattice limit, and only afterward increase the lists.  Under
centered sampling, \(\E_{\rm cen}r_H=e^{-D(H,p)}\), the convergence
of the pin interaction, and strict positivity of the Gaussian
\(r_H\) close the approximation to the full core pin in probability.
The restoration consequence transfers the exceptional events between
the centered topologies; H\"older with the uniform fixed-tilt
exponential moments transfers them to the sampled tilted laws.
This use of event transfer is consistent with the full density
\eqref{ent:tilted-guard-change}.  On exceptional data use
\(\delta_i\le1\).  Along an admissible sequence with
\(y_s^\ell\to\bar y^\ell\) in the compact lifted set, after taking
a fixed topology subsequence, this gives
\[
 \limsup_{s\to\infty}
   \E_{\nu_{t,y_s^\ell}^p}\delta_i^r
 \le \E_{\nu_{t,{\rm g},\bar y^\ell}^p}
                         (\Theta_{i,{\rm g}}^{\rm up})^r ,
\]
where the Gaussian upper envelope is
\[
 \Theta_{i,{\rm g}}^{\rm up}
   =\min\left(1,\ d_{i,{\rm g}}-\log r_H^{\rm g}
                                  +|\Lambda_i^{\rm g}|\right).
\]
The cutoff logarithm has limit zero.  This envelope uses the
limiting pin-interaction bound and does not require convergence of
\(c_{i,s}\).  With
\(\|y^\ell\|_2\le\operatorname{poly}(P)
                    (1+\sqrt{\mathcal E_b(w)})\), the preceding
prediction estimate bounds their norm by
\[
\begin{split}
 \|\Theta_{i,{\rm g}}^{\rm up}\|_{L^r}
 &\le Ce^{-cR}+\tfrac12\|Y_i\|_{L^{2r}}^2
       +C\|y^\ell\|_2\|Y_i\|_{L^r}
       +Ce^{-cR}\|y^\ell\|_2^2\\
 &\le CP^Cr^Ce^{-c_0R}(1+\mathcal E_b(w))^C .
\end{split}
\]
Compact subsequences and the uniform exponential moments make this
fixed-target conclusion uniform.  The slack \(\eta\) gives
\eqref{ent:one-mark-moment}; no rate in \(s\), limit of \(r_\Delta\)
or \(\delta_i\), or convergence uniform in all orders is used.

Choose \(\eta\le e^{-c_0R}\) at these fixed parameters and orders.
Let \(j=j_{\rm deep}(b)\) in an ordinary good halo box.  For each
mark, the current positive box topology is one member of its
adjacent pair with all other choices fixed.  Thus every factor is
estimated under a marginal of the same current box law in
\(Q_b^\tau\).  H\"older with order \(j\) when \(j\ge2\), and order
two when \(j=1\), gives
\begin{equation}\label{ent:good-box}
\begin{split}
 Q_b^\tau(y_b)
 &\le C\left[CP^C(j\vee2)^Ce^{-c_0R}
                    (1+\mathcal E_b(w))^C\right]^j\\
 &\le C\exp\{-cRj+R^{-12}\mathcal E_b(w)\}.
\end{split}
\end{equation}
For \(j=0\), the separate compact bound
\(\mathcal A_{G_b}^\tau\le C\) gives the same result.  No independence
of guards is used.  The last inequality follows by maximizing the
polynomial factor after multiplication by
\(e^{-R^{-12}\mathcal E_b}\): since \(j\le\operatorname{poly}(P)\),
the cost is \(O(j\log R)\), absorbed by a fixed fraction of
\(c_0Rj\).  This is the good-box argument of \cite[Section~9]{B}
with the smaller allowance in \cite[Equation~(5.20)]{C}.

For an isolated selected tile, \(\mathcal E_b\) controls its whole
observation vector modulo the component period because every original
tile face is present in that energy.  Its fixed-tile normalized factor
therefore tends uniformly to one on good data, giving
\(\lvert T_Q^\tau-1\rvert\le e^{-cR}\) for sufficiently large \(s\).
There are only finitely many shapes and mark lists at fixed
\(R,P,U_{\rm big}\), so all these bounds hold for every sufficiently
large admissible \(s\).  In a compulsory box we discard all
good-box gains.  Only a number depending on \(R\) of gains is lost.

\paragraph{Ordinary bad and compulsory boxes.}
If \(\mathcal E_b(w)>U_{\rm big}^2\), discard the conditional
rarity and use \eqref{ent:large-field-reserve} for the
positive factor of a clipped halo piece in
\eqref{ent:positive-topology}.  For an isolated selected tile \(Q\),
use \(\lvert T_Q^\tau-1\rvert\le T_Q^\tau+1\) and apply
\eqref{ent:large-field-reserve} on the intact graph \(Q\).
Its full-tile classification is retained on the constant-one branch.
On an ordinary bad box of either kind insert
\begin{equation}\label{ent:threshold-charge}
  1\le
     \exp\{\varepsilon_1(\mathcal E_b(w)-U_{\rm big}^2)\},
\end{equation}
where \(\varepsilon_1>0\) is fixed small compared with the
gap in \eqref{ent:reserve-order}.  Every compulsory box, regardless
of its energy, uses the broad bound
\eqref{ent:large-field-reserve} without
\eqref{ent:threshold-charge}; its determinant and lost
gains will be a fixed compulsory cost.  The intact isolated source
tile is always handled by this compulsory bound, without a threshold.
Thus the bad-box
coefficient is \(1-c_*<h_1\) everywhere a broad energy is
used.  This strict inequality is precisely where the
observation cutoff is needed.

\paragraph{The local quadratic cost and Gaussian averaging.}
Write the argument of \eqref{ent:averaged-interaction} as
\[
 y=y^0+B_sv_n,\qquad
 y^0=B_s\alpha(\psi+\zeta_{<m})+e,\qquad
 w=\nabla_{\rm c}y^0+
                 (\nabla_{\rm c}B_sv_n+\tau).
\]
For the centered sector the deterministic summand is zero.
First apply quadratic Cauchy--Schwarz to the shifted broad energies
and face squares,
using a fixed small fraction of the gap between \(1-c_*\) and
\(h_1\).  Reserve a further fraction for the face costs.  In the
following bound the cost in \(y^0\) is the coefficient after this
split.
For each fixed positive branch and box translate, the ordinary boxes'
good/bad predicates partition the real observation space into finitely
many classes; compulsory boxes retain their fixed status.  Keep each
class indicator through its class-dependent bounds above and the split
just made.  Fix one term of this partition, and let \(q_\pi(y^0)\) be exactly the sum of the
positive quadratic costs retained after this split: the centered
broad energies on the actual bad and compulsory graphs, the
good-box allowances, the localized proxy, seam, and link costs, and
the positive part of \eqref{ent:threshold-charge}.  Its coefficients
are those in the preceding bounds; the negative threshold constants
and deterministic source costs are kept outside \(q_\pi\).  Before the
Gaussian integral for this fixed \(q_\pi\), replace its class indicator
by one.  The finite class count is charged to the existing per-box
counting cost below.  This positive quadratic form satisfies
\begin{equation}\label{ent:dominating-open-energy}
             q_\pi(y^0)\le\tfrac12 h_1 E_D^0(y^0),
                  \qquad D=D_m(X).
\end{equation}
Here \(D\) is enlarged to whole cells and contains
\(D_{\mathrm{phys}}(\pi)\), the full link paths and collars, and every
support used by the quadratic costs.  There are
two form comparisons behind this statement.  Restrict a
trial field for \(E_D^0\) to all cut halo pieces receiving a bad or
compulsory bound, all intact selected tiles receiving a bad or
compulsory bound, and any compulsory isolated source tile.
Their actual graphs are disjoint: selected tiles are outside \(G\)
and are mutually isolated.  Thus their internal edges and observation
rows are counted at most once, and the sum of their minimized
energies is at most \(E_D^0\).  Comparison-only proxy and seam collars
can overlap, but they enter through the small face coefficients with
bounded per-face multiplicity.  On two adjacent underlying
cells, cell Poincar\'e and their observation residuals give
\[
 |y_j^0-y_i^0|^2\le C\left\{
   a^{-1}\sum_{e\text{ in their union}}|\nabla u(e)|^2
       +\sum_{k=i,j}|(B_su)_k-y_k^0|^2\right\}.
\]
The union in this display includes their common underlying face.
Sum this with the localized coefficients already bounded
per face and minimize over \(u\).  For large \(R\) and
small \(\varepsilon_1\), these costs fit into the strict
gap above the broad coefficient \(1-c_*\).  This proves
\eqref{ent:dominating-open-energy}.

For the remaining deterministic part of the split, trying \(v_n\)
in a cut energy gives
\[
 E_G^\tau(B_sv_n)\le a^{-1}\|\eta_n\|_G^2.
\]
Likewise
\((B_sv_n)_j-(B_sv_n)_i+\tau_f=\sum_e\Phi_f(e)\eta_n(e)\),
so \eqref{ent:cell-flow} bounds a deterministic face cost
by the energy on its two cells.  Away from the compulsory
neighborhood, \eqref{ent:unit-decay} makes this at most a
fixed constant per halo or selected-tile item and per block along a
link cover.  More precisely, an ordinary physical portion with \(O(R^2)\)
present cells at cell distance at least a fixed multiple
of \(P\) costs at most \(CR^2/P^2\).  The portions have
bounded overlap.  This uses the actual cells of a clipped halo or
an intact tile, without filling its containing box.
The exponentially small link coefficients give an even
smaller cost, or can be paid per block of their linear
covers.  In a pattern containing \(I_\star\), the deterministic energy
in the fixed source neighborhood containing its compulsory groups is at most
\(C_R\log(2+s)\), by summing \eqref{ent:unit-decay} from the source
to distance \(O_R(s)\).  It is charged to the labelled component.
For an unlabelled unit-sector pattern, compatibility with
\(\mathfrak C_{\star,m}\) puts every physical portion at cell distance
at least a fixed multiple of \(P\) from the source, so the preceding
ordinary per-item estimate applies to all its deterministic costs.

Write \(y^0(z)=b+\mathsf S z\), where \(b=B_s\alpha\psi\) and
\(z\) consists of standard Gaussian/Cameron coordinates on the
covariance support for
\((\zeta_{<m},e)\).  The linear map \(\mathsf S\) includes the
observation averaging and noise.  Define the positive matrix
\[
       \mathsf M_\pi=\nabla_z^2\!\left[p_1q_\pi(b+\mathsf S z)\right].
\]
It is the Hessian of the actual localized cost in the \(p_1\)
H\"older factor, rather than the Hessian of the dominating open
energy.  Form domination in \eqref{ent:dominating-open-energy}
and the covariance contraction of Lemma~\ref{ent:entry-regulator}
give
\[
                 0\le \mathsf M_\pi\le p_1h_1 I<I .
\]
The strict inequality follows from
\(p_1<p_2\) and \(h_1<h_*<1/p_2\).  For the current ordinary
regulator factor \(W\), the chosen H\"older inequality is
\[
 \E[We^{q_\pi}]
 \le
 \bigl(\E e^{p_1q_\pi}\bigr)^{1/p_1}
 \bigl(\E W^{p_1/(p_1-1)}\bigr)^{(p_1-1)/p_1}.
\]
Gaussian completion in its \(p_1\) factor is exactly
\begin{equation}\label{ent:entry-completion}
\begin{split}
 \left(\E_z e^{p_1q_\pi(b+\mathsf S z)}\right)^{1/p_1}
 ={}&\det(I-\mathsf M_\pi)^{-1/(2p_1)}\\
 &\quad{}\times
 \exp\left\{\sup_z\left[
       q_\pi(b+\mathsf S z)-\frac{\|z\|^2}{2p_1}\right]\right\}.
\end{split}
\end{equation}
The supremum is at most \(V_{j_0}(X,\psi)\), by
\eqref{eq:observation-regulator} and \(h_1<h_*\).  The other
H\"older factor is the ordinary regulator with its reserved
exponent.  Thus the completion uses the open energy only to bound
the shifted exponential; the determinant keeps \(\mathsf M_\pi\) and its
localized trace.
Its trace is bounded by
\begin{equation}\label{ent:entry-trace}
\begin{split}
 \tr \mathsf M_\pi\le {}&
 CP^2 N_{\rm bad}
 +CR^{-12}P^2N_{\rm good}
 +C(R/P)R^2(j_1+j_2)\\
 &+C\sum_\ell e^{-c(R+d_\ell)}
                         \operatorname{poly}(d_\ell)
 +C_R1_{\{\star\}} .
\end{split}
\end{equation}
Here \(N_{\rm bad}\) and \(N_{\rm good}\) count ordinary bad and good
clipped halo boxes and auxiliary tile boxes, using their respective
energies above; the compulsory boxes are in the last term.  Boxes met
solely by links contribute through the fourth term.  A tile contributes
one auxiliary box even on its constant-one branch.  Each broad halo
energy and its threshold charge have rank at most \(CP^2\) in cell
observations, and the contraction bounds every eigenvalue.
For an intact tile the rank is \(O(R^2)\), hence also at most \(CP^2\).
The good allowance has coefficient \(R^{-12}\) on at most
\(CP^2\) face contrasts.  Each such contrast has bounded variance
under the reference Gaussian and noise.  Indeed, for an underlying
face \(f=(i,j)\), the flow in Equation~\eqref{ent:cell-flow} and
\(A^{1/2}C_{<m}A^{1/2}\le I\) give
\[
 \Var\bigl((B_s\alpha\zeta_{<m})_j-(B_s\alpha\zeta_{<m})_i
                              +e_j-e_i\bigr)
       \le a\|\Phi_f\|_2^2+2\le C.
\]
The contrast annihilates constants.  Its underlying face belongs to
the full reference graph and the open support \(D\), even when the
positive topology has deleted that face.  The same bound applies to
the seam and link costs.  There are \(O(R^2)\) present
cells per halo for the seam trace.  At a fixed face, the sum of the
reserved coefficients over all possible links is bounded by
\(Ce^{-cR}\sum_d\operatorname{poly}(d)e^{-cd}\): only polynomially
many tagged descriptions of length \(d\) meet that face.  This controls
every selected subset of links, independently of its cardinality, and
gives the fourth trace term.  Each marked halo intersects a bounded
number of \(P\)-boxes and each selected tile adds one auxiliary box, so
\(N_{\rm good}\le C(j_1+j_2)+C_R1_{\{\star\}}\).
The actual good quadratic remains inside \(D\); the
\(P^2\) count is only an upper bound for its trace.
These are the trace charges of
\cite[Equation~(9.26)]{B} and
\cite[Equation~(5.24)]{C} for the same reference forms
in \(y^0\).

The unit junction term in \eqref{ent:face-decomposition}
does not enter \(\mathsf M_\pi\).  Equation~\eqref{ent:fixed-circulation}
makes it a constant in these standard Gaussian coordinates,
so its Hessian is zero even though its coefficient is
\(O(\log s)\).  A cut that changes its presence pays a
factor \(s^{C_R}\), not a determinant.  All deterministic source
costs were separated before \(q_\pi\) and remain outside this
completion.  The background \(b=B_s\alpha\psi\) supplies the affine
terms in \(z\) already included in the supremum.  In particular no determinant
in \eqref{ent:entry-trace} counts the microscopic primary
sites or the subdivision vertices.

For a positive matrix with norm at most \(1-\delta\),
\(-\log\det(I-\mathsf M_\pi)\le\delta^{-1}\tr \mathsf M_\pi\).  Thus
\eqref{ent:entry-trace}, with the fixed factor \(1/(2p_1)\) in
\eqref{ent:entry-completion}, bounds the logarithmic determinant.
A bad halo box meets at most a polynomial in \(P\) number of original
marks, and a bad auxiliary tile box carries one selected tile, including
its constant-one branch.  Choose the finite \(U_{\rm big}\), now and before
choosing \(s\), so that the negative constant
\(\varepsilon_1U_{\rm big}^2\) in
\eqref{ent:threshold-charge} pays the \(CP^2\)
determinant, the lost gains of those marks and tiles, and the per-box
counting cost, with a further \(e^{-cR}\) per such item.
The link gains remain in their separate trace and counting bounds.
Since \(P=R^5\),
both \(R^{-12}P^2\) and \((R/P)R^2\) are \(O(R^{-2})\).
They do not consume the remaining ordinary gains.  The
compulsory boxes pay their whole fixed determinant and
lost gains in \(C_R\), and the deterministic source costs
give \(s^{C_R}\).  We obtain the real estimate
\begin{equation}\label{ent:real-pattern}
 |F_\pi(\psi)|
 \le \mathcal C_\pi(s)\,
 C^{j_1+j_2}e^{-cR(j_1+j_2)}
 \prod_\ell Ce^{-c(R+d_\ell)}
 W_m^{\kappa_0}(X,\psi)e^{V_{j_0}(X,\psi)},
 \qquad \kappa_0<\kappa .
\end{equation}
There is strict reserve in both regulators.  In particular,
for a real direction \(f\) with \(\|f\|_{m,D}\le1\),
replacing the background by \(\psi+xf\) while keeping the
regulator evaluated at \(\psi\) costs at most
\begin{equation}\label{ent:real-translation}
                              e^{Cx^2|X|}.
\end{equation}
This follows by Cauchy--Schwarz in the optimized quadratic
form using \(h_1<h_*\); the direction energy is
\(\ip{f}{A_Df}\le C|X|\).  The ordinary regulator has
the same spare exponent.

\paragraph{Complex directions and the derivative norm.}
For real \(y,z\), the finite theta parity inequality on
the chain subdivision gives
\begin{equation}\label{ent:complex-comparison}
 \frac{|Z_G^\tau(y+iz)|}{Z_G^\tau(y)}
                 \le\frac{Z_G^0(iz)}{Z_G^0(0)}.
\end{equation}
To include the sector, restrict the affine lattice to its
parallel span and move its real displacement into the
real theta tilt.  The theta inequality in
\cite[Equation~(9.29)]{B} applies to that tilt.  A
displacement normal to the span contributes the same
imaginary Gaussian factor on the left and the centered
right, with a phase of modulus one.  This proves
\eqref{ent:complex-comparison} on the affine lattice.
Its dominated limit at the fixed primary graph proves
the Bessel version, as in
\cite[Equation~(5.26)]{C}.  Every freely translated
component is confined by the observations during this
comparison.

The centered imaginary ratio is its Gaussian prefactor
times a positive centered characteristic function.
Centered pinning increases that characteristic function,
by Proposition~\ref{prop:comparisons}.  Expand the
differences into their positive topologies and temporarily
discard the \(\delta_i\) rarity, while retaining the link
gains.  For the halo factors, pin whole strips of width \(R/4\) on translates
of the \(P\)-box walls and cut through the middle of
these zero strips, using only their intersections with the fixed
physical envelopes.  The strips and every auxiliary cut
are unions of whole observation cells, so no average
straddles a guard or a new component.  Only the primary heights in
the present strip cells are pinned; every observation row is retained.
On the resulting
fixed boxes, Proposition~\ref{prop:height-inputs} gives
the Gaussian imaginary ratio uniformly on bounded
imaginary arguments after \(s\) is sufficiently large.
For clarity, let \(S\) be the present strip cells in a halo graph
\(G\), and let \(E_{G,S}^0\) denote its Gaussian observation energy
with those primary cells pinned to zero.  The unpinned minimizer is
\(u_z=T_G^{-1}B_s^*z\), with \(T_G=A_G/a+B_s^*B_s\).
Choose a cell cutoff \(\chi\) equal to one on \(S\) and supported
in a fixed cell collar \(N_G(S)\) inside the present graph, with
microscopic slopes at most \(C/s\) on retained edges.
The trial \((1-\chi)u_z\) is pinned-admissible, and quadratic
completion gives
\[
 E_{G,S}^0(z)-E_G^0(z)
       \le \|\chi u_z\|_{T_G}^2
       \le C\|z\|_\infty^2\,|N_G(S)|.
\]
The last bound follows from cell coercivity and the exponential
locality of the columns \(T_G^{-1}B_s^*1_Q\), proved by the same
cell-cutoff argument as Lemma~\ref{ent:face-localization}.  It holds
on fragmented unions and at free walls because it uses only present
cells and retained edges.  Averaging \(|N_G(S)|\) over wall
translates therefore bounds the additional Gaussian wall energy by
\[
       C\|z\|_\infty^2(R/P)
       \#\{\text{halo observation cells in }D_{\mathrm{phys}}(\pi)\}.
\]
This is the imaginary-wall estimate in
\cite[Equation~(9.31)]{B}; it follows from massive
locality away from the strips and the bounded energy
per cell near them.  The count uses actual halo cells and their fixed
physical comparison regions, with no empty container cells; link-only
cells stay in the separate link bound.  The displayed count costs \(O_H(R^3/P)=O_H(R^{-2})\)
per halo when the imaginary direction is in a fixed
strip of height \(H\).

The Gaussian shapes cancel the imaginary Gaussian ratio.
In fact the quadratic continuation of their affine form obeys exactly
\begin{equation}\label{ent:imaginary-shape}
          \Re E_G^\tau(y+iz)=E_G^\tau(y)-E_G^0(z).
\end{equation}
The imaginary face data are \(\nabla_{\rm c}z\) and
have zero circulation, so the second term is the centered
one also in \eqref{ent:face-decomposition}.  The bounded
matrix errors are the small localized costs already
charged above.

For a halo topology, expanding the signed differences has removed
only the conditional \(\delta_i\) expectation.  After the theta
comparison and the Gaussian shape cancellation, its physical real
factor is still \(\mathcal A_G^\tau(y)\), multiplied by the
imaginary-wall and cutoff costs.  Apply the real seam comparison
\eqref{ent:seam-normalized} to that factor with the delta product
omitted.  Its per-box factor is then
\(\mathcal A_{G_b}^\tau(y_b)\) in place of \(Q_b^\tau(y_b)\).

For an intact isolated tile \(Q\), apply
\eqref{ent:complex-comparison} on \(Q\) itself.  Its centered imaginary
ratio has the fixed-tile Gaussian limit and cancels its Gaussian shape
by \eqref{ent:imaginary-shape}, without a \(P\)-wall strip.  The real
factor remains \(T_Q^\tau(y)\), so it uses exactly the full-tile good,
bad, or compulsory bound above; the constant-one branch has no imaginary ratio.
The same calculation applies to an intact source tile with its
compulsory bound.

On an ordinary good halo piece use the separate compact estimate
\(\mathcal A_{G_b}^\tau(y_b)\le C\), allowing the harmless factor
\(e^{R^{-12}\mathcal E_b(w)}\).  The mark gain \(e^{-cRj_{\rm deep}}\)
from \eqref{ent:good-box} is unavailable in this rough strip.
For an intact selected tile, the triangle bound keeps the real
terms \(T_Q^\tau(y)\) and \(1\) separate and likewise uses only
their compact bound on good data.  Only real bad and compulsory
boxes use the broad energy bound.  Thus their \(CP^2\) determinant
costs remain confined to those boxes, while ordinary good boxes
have at most the \(R^{-12}P^2\) trace cost in
\eqref{ent:entry-trace}.  The real rarity gains will be recovered
by interpolation from the real boundary.

The additional cutoff factors satisfy the exact elementary
bound
\[
            |\chi_s(x+it)|\le \chi_s(x)e^{t^2/2}.
\]
For a complex field translation by \((x+iy)f\), with
\(\|f\|_{m,D_m(X)}\le1\) and \(|y|\le H\), each imaginary
face difference is \(O_H(s/m)=O_H(R^{-1})\).
There are \(O(R^2)\) faces per halo or intact tile, so these extra
squares cost only \(O_H(1)\) per such item.  The compulsory
region has only a fixed additional cost depending on
\(R,H\).  Link costs retain their tagged-length gains.
We have therefore obtained the rough bound
\begin{equation}\label{ent:rough-strip}
\begin{split}
 |F_\pi(\psi+(x+iy)f)|
 \le {}& \mathcal C_\pi(s)
 \exp\{o(R)(j_1+j_2)+C_H(1+x^2)|X|\}\\
 &\quad{}\times
 \prod_\ell Ce^{-c(R+d_\ell)}
 W_m^\kappa(X,\psi)e^{V_{j_0}(X,\psi)},
                    \qquad |y|\le H.
\end{split}
\end{equation}
The \(o(R)\) is for fixed \(H\) as \(R\) increases; it
includes the preceding wall and finite per-halo or per-tile costs.
Finite-graph analyticity follows from absolute
summation with the large-field reserve.  The same bounds
justify analyticity after the Gaussian averaging and
the absolutely summable links.

Choose \(r_0>4e h_0\) and \(H>4r_0\) before choosing
\(R\).  Apply the three-lines theorem in the upper and
lower strips to
\[
           F_\pi(\psi+zf)\exp\{-\gamma_Hz^2|X|\},
                         \qquad \gamma_H>C_H.
\]
The real boundary uses \eqref{ent:real-pattern} and
\eqref{ent:real-translation}; the other boundary uses
\eqref{ent:rough-strip}.  On \(|z|\le r_0\) the harmonic
weight of the real boundary is at least \(3/4\).
The item gains therefore survive with a smaller \(c\).
Removing the damping costs \(e^{C|X|}\), which is
absorbed by the gains and \(\mathcal C_\pi(s)\), since
\[
 |X|\le C(j_1+j_2)+C\sum_\ell(1+d_\ell/R)
                         +C_R1_{\{\star\}}.
\]
Cauchy's inequality bounds the diagonal derivative of
order \(k\) by \(k!r_0^{-k}\) times this recovered
pattern bound.  Real polarization costs a further
\(k^k/k!\).  Its contribution to the point norm is thus
at most the pattern bound times
\((h_0/r_0)^k k^k/k!\), whose sum converges because
\(r_0>4e h_0\).  This proves
\eqref{ent:analytic-pattern-bound}.
\end{proof}

\begin{proof}[Completion of Theorem~\ref{ent:initialization}]
It remains to sum connected patterns and check the norm cost
of the freely filled translated entries.  Assign to each
ordinary halo item and link the positive weight given by
Lemma~\ref{ent:pattern-estimate}, including its fixed
analytic and geometric constants.  For every entry block
\(B\), their placement sum satisfies
\begin{equation}\label{ent:placement}
 \sum_{\substack{I:\ \operatorname{cov}_m(I)\text{ meets}\\
                 \text{the connection neighborhood of }B}}
 w(I)(2A_w)^{|\operatorname{cov}_m(I)|}
              e^{|\operatorname{cov}_m(I)|}
                       \le e^{-cR}=:\epsilon_R .
\end{equation}
A halo has bounded entry-block diameter and only polynomially
many observation-cell roots and event labels in a block.
A link of tagged length \(d\) has
\(O(1+d/R)\) covered blocks and only polynomially many
endpoint and path choices at that length, including its
period images or wall tag.  Their weights and cover factors
are paid by \(e^{-c(R+d)}\).  These are counts in observation
cells, independent of \(s\).  This proves
\eqref{ent:placement} by the connected-placement estimate
of \cite[Equation~(9.34)]{B}.

For clarity, root an ordinary connected collection at an
item meeting a fixed block and choose a canonical spanning
tree of its item adjacency graph.  If the sum of descendant
trees attached at one block is \(T\), the unordered
children of an item with cover \(S\) cost at most
\(e^{|S|T}\).  Starting with depth zero,
\eqref{ent:placement} and its reserved \(e^{|S|}\)
show inductively that \(T\le\epsilon_R<1\) at every
depth.  Monotone convergence bounds all finite trees.
The factor \(2^{|S|}\) pays the regrouping into occupied
subsets.  This gives the norm bound \(Ce^{-cR}\)
for the ordinary connected sum, uniformly on plane,
torus, and square covers.

For a labelled collection the root is the single
compulsory cover.  All of its \(|\mathfrak C_{\star,m}|
\le C_R\) blocks remain occupied.  They supply at most
\(C_R\) choices for child attachment and a fixed
weight cost depending on \(R\); the same descendant
tree bound sums every attached ordinary item.  Multiplying
by \(\mathcal C_\pi(s)=C_Rs^{C_R}\) proves
\eqref{ent:label-norm}.  In particular this summation
never drops the padded or tagged part of
\eqref{ent:occupied-label-cover}.  It also proves absolute
integrated convergence of \eqref{ent:face-links}.  On a
finite graph one may first sum the finite positive
topologies, then integrate the noise and high Gaussian,
and finally sum the links.  Tonelli for absolute values
and the proved bound justify all exchanges in the signed
identity.

Finally consider an entry of \(K^\tau\) that was filled
by the translated prescription because it cannot
coexist with the compulsory cover.  Put \(D=D_m(X)\), and let
\(F_X\) be the actual field and difference-stencil support of
\(K_{\mathcal G,m}^0(X,\cdot)\).  By the local calculus and the full-cover
assignment, \(F_X\) lies in the sites of \(X\) plus a fixed microscopic
collar of width at most \(\rho m\), with a fixed \(\rho<1\).
Since the rectangle spanned by \(D\) misses the source, choose a
signed coordinate \(t\) for which that rectangle lies in \(t>0\).
The two padding layers on the source-facing side each have width at
least \(m\), including at irregular end blocks.  They cannot be
clipped by the source-facing free wall: reaching that wall would make
the padded rectangle cross the interior source in this coordinate.
Consequently
\[
             F_X\subset\{t\ge(2-\rho)m-C\}\subset\{t\ge m\}
\]
for all sufficiently large admissible \(m\).  This separates the actual
field inputs from the source even when the outer padding \(D\) approaches it.

Choose the physical lift \(\lambda_X\) on the source-free half-rectangle
\(\{t\ge m/3\}\), with its additive period aligned to the lift in
\eqref{ent:translated-copy} on \(F_X\).  Lemma~\ref{ent:center-background}
gives uniformly bounded oscillation there, including the possible
one-period unwrapping across the fixed cut.  Choose a constant \(c\) so
that \(|\lambda_X-c|\le C\) on this half-rectangle.  Let
\(\chi\in C^2(\mathbb R)\) be fixed, equal to zero on
\((-\infty,1/2]\) and to one on \([1,\infty)\).  On the lattice define
\[
 \lambda_{\rm ext}
       =c+\chi(t/m)(\lambda_X-c)
\]
where \(t\ge m/3\), and continue it as \(c\) below that half-rectangle.
The formulas agree on a neighborhood of their join because \(\chi=0\)
through \(t=m/2\).  On \(F_X\), \(\lambda_{\rm ext}=\lambda_X\).
Locality therefore gives
\[
 K_{\mathcal G,m}^0(X,\psi+\lambda_{\rm ext})
       =K_{\mathcal G,m}^0(X,\psi+\lambda_X)
\]
as functions of \(\psi\), and hence gives equality of all derivatives
in \eqref{eq:point-norm}.  The exact copy identity retains its physical
lift \(\lambda_X\); \(\lambda_{\rm ext}\) is a representative for
estimating the same activity.

The bounds in Lemma~\ref{ent:center-background} on \(t\ge m/3\),
the bounded oscillation, and the discrete product rules for the fixed
cutoff give
\[
       |\nabla\lambda_{\rm ext}|\le C/m,\qquad
       |\nabla^2\lambda_{\rm ext}|\le C/m^2
\]
throughout \(D\) and every boundary and stencil neighborhood used by
the norm.  At a free wall normal to \(t\), the cutoff is constant in
its stencil neighborhood under the volume hypothesis; at a perpendicular
wall it is independent of the normal coordinate.  The even reflected
convention thus preserves these bounds.  There are \(O(m^2|X|)\)
sites in \(D\), \(O(m|X|)\) boundary edges, and \(O(|X|)\) blocks.
For the quantities in the ordinary regulator of \cite[Equation~(4.10)]{B},
these counts give
\[
 g_D(\lambda_{\rm ext})+P_m(D,\lambda_{\rm ext})
       +T_m(D,\lambda_{\rm ext})\le C|X|.
\]
Trying \(\alpha\lambda_{\rm ext}\) in the open observation minimum,
with zero observation residual, also gives
\[
                  E_D^0(B_s\alpha\lambda_{\rm ext})\le C|X|.
\]

Write \(V_m^{(h)}\) for \eqref{eq:observation-regulator} with coefficient
\(h\).  Quadratic Cauchy--Schwarz and the optimized form of \(V\) in
\cite[Equation~(9.16)]{B}, using its open-energy bound from
\cite[Lemma~9.2]{B}, now give
\[
\begin{split}
 W_m^\kappa(X,\psi+\lambda_{\rm ext})
       &\le e^{C|X|}W_m^{\kappa'}(X,\psi),\\
 e^{V_m^{(h_*)}(X,\psi+\lambda_{\rm ext})}
       &\le e^{C|X|}e^{V_m^{(h_*')}(X,\psi)}
\end{split}
\]
for the reserved admissible \(\kappa'>\kappa\) and \(h_*'>h_*\).
Common values do not affect either regulator.  The stronger ordinary weight, fixed before
\(R\), pays \(e^{C|X|}\) on passing to the sector
weight.  This proves the small norm bound also on
those filled entries.

The scalar and subset identities, the fully occupied
label support, and the copy identities were established
by the finite construction.  The estimates now justify
their integration and give all claimed norms.  Increasing
the fixed \(R\) makes \(Ce^{-cR}\) smaller than any
of the prescribed activity balls.  The finite threshold
is then chosen, and the fixed-geometry limits make
the proof valid for every sufficiently large admissible
\(s\).  This completes the theorem.
\end{proof}
